\documentclass[11pt,a4paper]{article}

\usepackage[margin=1in]{geometry}
\usepackage[T1]{fontenc}
\usepackage[utf8]{inputenc}
\usepackage{lmodern}
\usepackage{microtype}
\usepackage{amsmath,amssymb,amsthm,mathtools}
\usepackage[most]{empheq}
\usepackage{booktabs}
\usepackage{enumitem}
\usepackage{xcolor}
\usepackage{framed}
\usepackage{graphicx}
\usepackage{float}
\usepackage[hypcap=false]{caption}
\usepackage{hyperref}
\usepackage{titlesec}

\hypersetup{
  colorlinks=true,
  linkcolor=blue!50!black,
  urlcolor=blue!60!black,
  pdftitle={Core publishable work: BDC and burst-aware BMVR},
  pdfauthor={Adam Aldridge}
}

% ---------------------------------------------------------------------------
% Macros
% ---------------------------------------------------------------------------
\newcommand{\dd}{{\rm d}}
\newcommand{\ee}{{\rm e}}
\newcommand{\E}{\mathbb{E}}
\newcommand{\Prob}[1]{\mathbb{P}\!\left\{#1\right\}}
\newcommand{\xt}{X_t}
\newcommand{\wt}{W_t}
\newcommand{\Ihat}{\widehat{I}}
\newcommand{\Kb}{\mathcal{K}}
\newcommand{\Icell}{\mathcal{I}}
\newcommand{\Vfree}{\mathcal{V}}
\newcommand{\Lap}[1]{\widetilde{#1}}
\newcommand{\ddt}[1]{\frac{\dd #1}{\dd t}}
\newcommand{\Fhyp}{{}_2F_1}
\newcommand{\Aa}{A}
\newcommand{\Bb}{B}
\newcommand{\highlight}[1]{\colorbox{blue!8}{\parbox{0.96\linewidth}{#1}}}

\theoremstyle{plain}
\newtheorem{result}{Result}[section]
\theoremstyle{definition}
\newtheorem{definition}[result]{Definition}
\theoremstyle{remark}
\newtheorem{remark}[result]{Remark}

\setlength{\parindent}{0pt}
\setlength{\parskip}{0.55em}
\setlist[enumerate]{leftmargin=1.6em,itemsep=0.25em}
\setlist[itemize]{leftmargin=1.6em,itemsep=0.2em}

\title{\textbf{Core Publishable Work}\\[0.35em]
\large From the single-cell birth--death--catastrophe process\\
to a burst-aware Basic Model of Viral Replication}
\author{Summary drawn from Chapters~3--4 and the BMVR extension notes}
\date{\today}

\begin{document}
\maketitle
\vspace{-1.5em}

\begin{abstract}
\noindent
This document is a clean, self-contained summary of the results that are ready
(or nearly ready) for publication. It has two layers.

\begin{enumerate}[leftmargin=1.4em]
\item \textbf{Single-cell BDC theory.} Closed forms for survival, moments,
burst time, burst size, and the quasi-stationary distribution --- with several
simplifications, one major correction, and one striking identity that the thesis
chapters contain the ingredients for but never state.
\item \textbf{Population (BMVR) extension.} A renewal formulation that embeds
those single-cell kernels into a within-host model, with exact effective
parameters, an invariant basic reproduction number, and a sharp criterion for when
bursting beats constant budding at matched~$R_0$ (the ``flooding advantage'').
\end{enumerate}
Everything stated as a Result has been checked algebraically and numerically.
Open technical tasks and thesis errata are deliberately omitted here.
\end{abstract}

\tableofcontents
\clearpage

% ===========================================================================
\section{The story in one page}
% ===========================================================================

Many within-host models treat infected cells as black boxes that release virions
at a constant rate~$p$ (the Basic Model of Viral Replication, BMVR). Real
lytic pathogens --- \textit{Yersinia pestis}, \textit{Francisella tularensis},
lytic phage --- do something different: they grow inside a cell and dump their
load in a single burst. The question is whether that strategy can be modelled
exactly, and what it buys the pathogen early on, when stochastic extinction is
still likely.

The route taken here is:
\begin{enumerate}
\item model one infected cell as a classical \textbf{birth--death--catastrophe
(BDC)} process (birth rate~$\beta$, death~$\mu$, load-proportional burst
rate~$\delta$);
\item derive closed forms for everything that cell does --- survival, load,
release, burst timing and size;
\item feed those formulas into a \textbf{renewal (age-structured) BMVR}, so the
population model is no longer phenomenological but is forced by the
intracellular process;
\item read off what changes: effective release rate, $R_0$, growth rate, and
especially the probability that the infection dies out in the first generation.
\end{enumerate}

The headline biological result is simple to state:
\begin{center}
\fbox{\parbox{0.9\linewidth}{\centering
At fixed total output per cell, bursting lowers extinction probability relative
to constant budding \emph{if and only if}
\[
\frac{1}{\delta}\;<\;\frac{1}{\beta}+\frac{1}{\mu}.
\]
When $\mu=0$, bursting always wins.}}
\end{center}

Everything else in this document either supports that claim, or shows how the
same single-cell formulas yield effective BMVR parameters and testable
single-cell predictions.

% ===========================================================================
\section{Notation (used consistently throughout)}
% ===========================================================================

{\small
\renewcommand{\arraystretch}{1.2}
\begin{tabular}{@{}llp{0.42\textwidth}@{}}
\toprule
Symbol & Meaning & Notes \\
\midrule
$\beta,\mu,\delta$ & birth, death, catastrophe rates & per capita, intracellular \\
$\xt$ & intracellular particle count & CTMC \\
$R$ & absorbing rupture state & \\
$\wt$ & released count from one cell & $0$ until burst \\
$\tau$ & burst time & may be $\infty$ \\
$\Kb$ & burst size (random) & $=\xt$ at burst \\
$I(t)$ & no-burst probability & $\to b$ as $t\to\infty$ \\
$D(t)$ & cleared, no burst by $t$ & $=p_0(t)$ \\
$\Ihat(t)$ & still productively infected & $=I-D$ \\
$J(t),K(t)$ & $\E[\xt]$, $\E[\xt^2]$ & mean load moments \\
$V(t)$ & $\E[\wt]$ & cumulative mean release \\
$a,b$ & characteristic roots & $b<1<a$ \\
$\Aa,\Bb$ & $a-1$, $1-b$ & simplify almost everything \\
$\theta$ & $\beta(a-b)$ & growth exponent \\
$T$ & target-cell count (held fixed) & BMVR \\
$\gamma$ & infection rate & BMVR (not $\beta$) \\
$\Icell,\Vfree$ & infected cells, free virions & BMVR counts \\
$c$ & virion clearance & BMVR \\
\bottomrule
\end{tabular}
}

% ===========================================================================
\section{The single-cell process}
% ===========================================================================

\subsection{Definition}

Starting from $\xt=1$ (one founder particle), the process jumps
\[
n \;\to\; n+1 \quad\text{at rate }\beta n,\qquad
n \;\to\; n-1 \quad\text{at rate }\mu n,\qquad
n \;\to\; R \quad\text{at rate }\delta n.
\]
On the jump to~$R$ the cell bursts: $\wt$ is set to the pre-burst load and stays
there. The process is classical (Brockwell--Gani--Resnick; Karlin--Tavar\'e);
what matters for applications is the family of closed forms below.

\subsection{Root algebra}

\begin{equation}
\eta=\frac{\beta+\mu+\delta}{2\beta},\qquad
a=\eta+\sqrt{\eta^2-\tfrac{\mu}{\beta}},\qquad
b=\eta-\sqrt{\eta^2-\tfrac{\mu}{\beta}}.
\end{equation}
Write $\Aa:=a-1>0$, $\Bb:=1-b>0$, $\theta:=\beta(a-b)$. The identities that do
most of the work are
\begin{equation}
\label{eq:AB}
ab=\frac{\mu}{\beta},\qquad
\Aa\Bb=\frac{\delta}{\beta},\qquad
\Aa+\Bb=a-b.
\end{equation}
In particular $b<1<a$ always, and $\beta(1-a)(1-b)=-\delta$.

\subsection{Three fixation functions}

There are three natural absorbing-behaviour probabilities, not two:
\begin{align}
I(t)&=\Prob{\text{no burst by }t},\\
D(t)&=\Prob{\text{internal extinction by }t,\ \text{no burst}},\\
\Ihat(t)&=\Prob{\xt\notin\{0,R\}}=I(t)-D(t).
\end{align}
Both $I$ and $D$ solve the Riccati equation
\[
f'=\beta(f-a)(f-b)=\mu-(\beta+\mu+\delta)f+\beta f^2,
\]
with $I(0)=1$ and $D(0)=0$. Writing $w:=\ee^{\theta t}\ge 1$:
\begin{result}[Closed forms for $I$, $D$, $\Ihat$]
\begin{equation}
I(t)=\frac{a\Bb+b\Aa w}{\Bb+\Aa w},\qquad
D(t)=\frac{ab\,(w-1)}{aw-b},\qquad
\Ihat(t)=\frac{(a-b)^2\,w}{(\Bb+\Aa w)(aw-b)}.
\end{equation}
Limits: $I(\infty)=D(\infty)=b$, $\Ihat(\infty)=0$, $\Ihat(0)=1$.
\end{result}

\begin{remark}[Correction of thesis formula]
The formula $\Ihat(t)=\hat\eta\big/\bigl(1-(1-\hat\eta)\ee^{\hat\eta\beta t}\bigr)$
used in Chapter~3 is wrong whenever $\mu>0$. It comes from treating $\Ihat$ as if
it obeyed the same autonomous Riccati as~$I$. The correct branching term on a
birth is $I^2-D^2$, not $\Ihat^2$. The error is invisible when $\mu=0$ (then
$D\equiv 0$). For $X_0=k$ founders one has $\Ihat_k=I^k-D^k$, not $\Ihat^{\,k}$.
\end{remark}

\subsection{Moments and mean release}

\begin{result}[First and second moments; mean release]
\begin{align}
J(t)&=\E[\xt]=\frac{(a-b)^2\,w}{(\Bb+\Aa w)^2}
=-\frac{\beta}{\delta}(I-a)(I-b),\\[0.3em]
K(t)&=\E[\xt^2]=\Bigl[1+\tfrac{2\beta}{\delta}(1-I)\Bigr]J,\\[0.3em]
V(t)&=\E[\wt]=(1-I)\Bigl[1+\tfrac{\beta}{\delta}(1-I)\Bigr].
\end{align}
\end{result}
Two useful simplifications: $J$ no longer carries the cluttered prefactor of the
thesis form (the identity $\Aa\Bb=\delta/\beta$ cancels it), and $V$ is a
quadratic in $(1-I)$ alone --- no separate $J$ or $\mu$.

\begin{result}[Lifetime-time release and conditional mean burst size]
\begin{equation}
V_\infty=\frac{a(1-b)}{a-1}=\frac{a\Bb}{\Aa},\qquad
\E[\Kb\mid\text{burst}]=\frac{V_\infty}{1-b}=\frac{a}{a-1}.
\end{equation}
$V_\infty$ includes realisations that clear internally and release nothing;
$\E[\Kb\mid\text{burst}]$ conditions them out.
\end{result}

% ===========================================================================
\section{Burst time, burst size, and the quasi-stationary distribution}
% ===========================================================================

\subsection{Burst time}

\begin{result}
$\Prob{\tau>t}=I(t)$, so the (defective) density of burst time is
\[
\varphi(t)=-I'(t)=\delta J(t),\qquad \int_0^\infty\varphi(t)\,\dd t=1-b.
\]
Mass $b$ sits at $\tau=\infty$ (internal clearance without burst).
\end{result}

\subsection{State probabilities are geometric at every time}

For $n\ge 1$,
\begin{equation}
p_n(t)=\Bigl(\frac{a-b}{b-aw}\Bigr)^{\!2}w\;P(t)^{\,n-1},
\qquad P(t)=\frac{1-w}{b-aw}.
\end{equation}
Conditional on $\xt\ge 1$, the load is geometric with ratio $P(t)$ at every~$t$,
and $P(t)\to 1/a$.

\subsection{The identity}

\begin{result}[Burst size $=$ quasi-stationary distribution]
Both are geometric on $\{1,2,\dots\}$ with ratio $1/a$:
\begin{align}
\mathrm{QSD}(n)
&=\lim_{t\to\infty}\Prob{\xt=n\mid \xt\notin\{0,R\}}
=(a-1)\,a^{-n},\\
\Prob{\Kb=k\mid\text{burst}}
&=(a-1)\,a^{-k}.
\end{align}
Unconditionally (including no-burst realisations),
$\Prob{\Kb=k}=(\delta/\beta)\,a^{-k}$ for $k\ge 1$, with mass $b$ at~$0$.
\end{result}

This is the cleanest single-cell result in the project. It is not obvious:
bursts are size-biased draws from the transient load, integrated against the
burst intensity $\delta n$. That the outcome is exactly the QSD deserves a
conceptual proof; the algebra is already complete (see below).

Immediate QSD moments:
\begin{equation}
\langle X\rangle_{\rm QS}=\frac{a}{a-1},\qquad
\mathrm{Var}_{\rm QS}(X)=\frac{a}{(a-1)^2}.
\end{equation}
Note $\langle X\rangle_{\rm QS}=\E[\Kb\mid\text{burst}]$. When $\mu=0$ both also
equal $V_\infty$; when $\mu>0$ they differ
(e.g.\ at $(\beta,\mu,\delta)=(1,0.2,0.05)$: $17.23$ versus $13.99$).

\subsection{Proof of the burst-size law (no simulation required)}

With $P(t)=(1-w)/(b-aw)$ one has
\[
\frac{\dd}{\dd t}\Bigl[\frac{P^k}{k\beta}\Bigr]=p_k(t).
\]
Hence
\[
\Prob{\Kb=k}
=\delta k\int_0^\infty p_k(t)\,\dd t
=\frac{\delta}{\beta}\bigl[P^k\bigr]_0^\infty
=\frac{\delta}{\beta}\,a^{-k}.
\]
Two independent checks close the loop:
\begin{enumerate}
\item $\sum_{k\ge 1}(\delta/\beta)a^{-k}=1-b$ (eventual-burst probability);
\item $\sum_{k\ge 1}k(\delta/\beta)a^{-k}=V_\infty$ (mean release).
\end{enumerate}

\subsection{Size-biasing and late bursts}

A cell bursts at rate $\delta\xt$, so given burst at time~$t$ the size is
size-biased:
\[
\E[\Kb\mid\tau=t]=\frac{K(t)}{J(t)}
\;\xrightarrow[t\to\infty]{}\;
1+\frac{2\beta}{\delta}(1-b).
\]
A very late burst is roughly twice the size of an average one --- the signature
of size-biasing a geometric.

% ===========================================================================
\section{From one cell to a population: the renewal BMVR}
% ===========================================================================

\subsection{Why a constant~$p$ substitution fails}

Chapter~4 attempted to replace the BMVR release rate~$p$ by the QSD mean
$\langle X\rangle_{\rm QS}$. That cannot work: $p$ has units of particles per
unit time, while $\langle X\rangle_{\rm QS}$ is a pure count. When $\mu=0$ one
even has $\langle X\rangle_{\rm QS}=V_\infty$, so the proposal substitutes total
lifetime output for an output \emph{rate}.

The deeper issue is that infected cells are not synchronised: at any host time
$t$ they have a distribution of ages (times since infection). The right object
is therefore an age-structured, or renewal, formulation.

\subsection{Classical BMVR (for comparison)}

In the simplified (target-cell-constant) Basic Model of Viral Replication one
has incidence $i(t)=\gamma T\,\Vfree(t)$ and the ODE system
\begin{equation}
\label{eq:BMVR-classical}
\begin{aligned}
\ddt{\Icell}
&= \gamma T\,\Vfree - d_{\Icell}\,\Icell,\\[0.25em]
\ddt{\Vfree}
&= p\,\Icell - c\,\Vfree.
\end{aligned}
\end{equation}
Here $p$ is a constant per-cell release rate and $d_{\Icell}$ a constant infected-cell
removal rate. Both are phenomenological: they are not derived from intracellular
dynamics.

\subsection{The two kernels from the BDC}

Both are already closed-form from the single-cell theory:
\begin{align}
\label{eq:kernels}
\Ihat(a)
&=\Prob{\text{cell of age }a\text{ is still productively infected}}
=\frac{(a-b)^2\,w}{(\Bb+\Aa w)(aw-b)},
\quad w=\ee^{\theta a},\\[0.3em]
g(a)
&=\delta K(a)=V'(a)
=\text{expected release flux at age }a.
\end{align}
The release kernel is correct because a cell bursts at rate $\delta\xt$ releasing
$\xt$ particles, so expected release in $[a,a+\dd a]$ is
$\delta\E[\xt^2]\,\dd a=\delta K(a)\,\dd a$. Explicitly,
\begin{equation}
K(a)=\Bigl[1+\tfrac{2\beta}{\delta}\bigl(1-I(a)\bigr)\Bigr]J(a),
\qquad
J(a)=\frac{(a-b)^2\,w}{(\Bb+\Aa w)^2}.
\end{equation}

\subsection{The new viral dynamics equations}

With target cells $T$ held constant, infection rate $\gamma$, virion clearance
$c$, and incidence
\begin{equation}
i(t)=\gamma T\,\Vfree(t),
\end{equation}
the burst-aware (renewal) model is the following integro-differential system.

\begin{result}[New viral dynamics --- burst-aware BMVR]
\label{res:new-vde}
\begin{empheq}[box=\fbox]{align}
\Icell(t)
&= \Icell_0\,\Ihat(t)
+\int_0^t i(t-\alpha)\,\Ihat(\alpha)\,\dd\alpha,
\label{eq:Icell-full}\\[0.5em]
\ddt{\Vfree}
&= \Icell_0\,g(t)
+\int_0^t i(t-\alpha)\,g(\alpha)\,\dd\alpha
-c\,\Vfree(t),
\label{eq:Vfree-full}
\end{empheq}
with kernels \eqref{eq:kernels} and $g=\delta K$. Equivalently, writing the
convolutions more compactly and using $i=\gamma T\,\Vfree$,
\begin{empheq}[box=\fbox]{align}
\Icell(t)
&= \Icell_0\,\Ihat(t)
+\bigl(\gamma T\,\Vfree * \Ihat\bigr)(t),
\label{eq:Icell-conv}\\[0.35em]
\ddt{\Vfree}
&= \Icell_0\,g(t)
+\bigl(\gamma T\,\Vfree * g\bigr)(t)
-c\,\Vfree(t).
\label{eq:Vfree-conv}
\end{empheq}
\end{result}

Here $\Icell_0$ is the number of cells already infected at $t=0$ (the initial
cohort); their contribution is $\Icell_0\Ihat(t)$ to the infected-cell count and
$\Icell_0 g(t)$ to the release flux. If the infection is seeded only by free
virions with $\Icell_0=0$, those terms drop and one recovers the pure convolution
form used in the working notes.

\paragraph{What each equation replaces.}
\begin{center}
\renewcommand{\arraystretch}{1.35}
\begin{tabular}{@{}p{0.42\textwidth}p{0.48\textwidth}@{}}
\toprule
Classical BMVR \eqref{eq:BMVR-classical} & New system \eqref{eq:Icell-full}--\eqref{eq:Vfree-full} \\
\midrule
$\Icell'=i-d_{\Icell}\Icell$
  (constant removal rate $d_{\Icell}$)
& $\Icell = \Icell_0\Ihat + i*\Ihat$
  (age-structured survival via $\Ihat$) \\[0.6em]
$\Vfree'=p\Icell-c\Vfree$
  (constant release rate $p$)
& $\Vfree' = \Icell_0 g + i*g - c\Vfree$
  (age-structured release via $g=\delta K$) \\
\bottomrule
\end{tabular}
\end{center}
Both constant BMVR terms are replaced: not only release $p\Icell$, but also cell
removal $d_{\Icell}\Icell$. Chapter~4's original attempt only ever proposed to
change the first.

\paragraph{Why the equations look different.}
In the classical model, infected cells are memoryless: every cell, regardless of
age, contributes the same $p$ and dies at the same $d_{\Icell}$. In the new model
a cell of age~$\alpha$ is still productively infected only with probability
$\Ihat(\alpha)$, and contributes expected release flux $g(\alpha)$. Because cells
are infected at different times, the population totals are convolutions of
incidence against these kernels --- not products of a stock $\Icell(t)$ with a
constant.

\paragraph{Optional differential form for $\Icell$.}
Differentiating \eqref{eq:Icell-full} (fundamental theorem + Leibniz rule) gives
\begin{equation}
\label{eq:Icell-diff}
\ddt{\Icell}
= i(t)
+\Icell_0\,\Ihat'(t)
+\int_0^t i(t-\alpha)\,\Ihat'(\alpha)\,\dd\alpha,
\end{equation}
so the incidence $i(t)$ appears as a source term exactly as in the classical
model, while the remaining terms encode age-dependent loss of productivity
(burst or internal clearance). There is in general no constant $d_{\Icell}$ for
which \eqref{eq:Icell-diff} reduces to $\Icell'=i-d_{\Icell}\Icell$ for all
histories of~$i$; such a reduction holds only for pure exponential solutions
(Section~\ref{sec:peff}).

\begin{remark}[What is and is not novel]
The renewal/age-structured framework itself is standard (McKendrick--von
Foerster; Nelson--Gilchrist--Perelson-type models). The novelty is that the
kernels $\Ihat$ and $g$ are derived in closed form from a mechanistic
intracellular stochastic process, rather than posited phenomenologically ---
together with the consequences below.
\end{remark}

% ===========================================================================
\section{Numerical overlays of the two systems}
\label{sec:overlays}
% ===========================================================================

To make the difference between \eqref{eq:BMVR-classical} and
\eqref{eq:Icell-full}--\eqref{eq:Vfree-full} concrete, both systems are
integrated side by side for a range of intracellular and population
parameters.

\paragraph{How to read every plot.}
\begin{center}
\renewcommand{\arraystretch}{1.35}
\begin{tabular}{@{}ll@{}}
\toprule
Style & Model \\
\midrule
\textbf{Solid navy line} & \textbf{NEW model} --- renewal / burst-aware BMVR
  (kernels $\Ihat$, $g=\delta K$) \\
\textbf{Dashed orange line with open circles} & \textbf{CLASSICAL BMVR} ---
  constant release rate $p$ and removal $d_{\Icell}$ \\
\bottomrule
\end{tabular}
\end{center}
Each panel also carries its own legend and the words ``NEW'' / ``CLASSICAL''
printed near the curves. The CLASSICAL parameters are the fairest
constant-parameter match at the $r=0$ end of the effective-rate map,
\begin{equation}
p = p_{\rm eff}(0)=\frac{\beta V_\infty}{\log\bigl(a/(a-1)\bigr)},
\qquad
d_{\Icell}=d_{\Icell,{\rm eff}}(0)=\frac{\beta}{\log\bigl(a/(a-1)\bigr)},
\end{equation}
so total lifetime output and mean productive lifetime agree. Figures are
generated by \texttt{figures/plot\_bmvr\_comparison.py}.

\subsection{Single-cell kernels (NEW model only)}

\begin{figure}[H]
\centering
\includegraphics[width=\textwidth]{figures/kernels.pdf}
\caption{Cell-survival kernel $\Ihat(a)$ and release kernel $g(a)=\delta K(a)$
for four intracellular parameter sets. These age-dependent kernels drive the
\textbf{NEW} model. The \textbf{CLASSICAL} model has no kernels --- only the
constants $p$ and $d_{\Icell}$.}
\label{fig:kernels}
\end{figure}

\subsection{Free virions and infected cells}

\begin{figure}[H]
\centering
\includegraphics[width=\textwidth]{figures/overlay_V.pdf}
\caption{Free-virion trajectories $\Vfree(t)$.
\textbf{Solid navy = NEW model} (renewal / burst-aware);
\textbf{dashed orange with circles = CLASSICAL BMVR} (constant $p$, $d$).
Six scenarios: supercritical growth, subcritical decay, fast burst, high
intracellular death, and free-virion-only seeding.}
\label{fig:overlay-V}
\end{figure}

\begin{figure}[H]
\centering
\includegraphics[width=\textwidth]{figures/overlay_I.pdf}
\caption{Infected-cell trajectories $\Icell(t)$ for the same six scenarios.
Same line code as Figure~\ref{fig:overlay-V}:
\textbf{solid navy = NEW}, \textbf{dashed orange = CLASSICAL}.}
\label{fig:overlay-I}
\end{figure}

\begin{figure}[H]
\centering
\includegraphics[width=\textwidth]{figures/overlay_V_with_naive.pdf}
\caption{Optional third curve (dotted green): CLASSICAL BMVR with the
dimensionally inconsistent Chapter~4 choice $p=\langle X\rangle_{\rm QS}$.
Still read as: \textbf{solid navy = NEW};
\textbf{dashed orange = CLASSICAL matched}; dotted green = CLASSICAL naive.}
\label{fig:overlay-V-naive}
\end{figure}

\subsection{Where matching fails}

Even with $p$ and $d_{\Icell}$ chosen to match mean lifetime output, the
CLASSICAL ODEs cannot reproduce the NEW-model transients: early on, young cells
dominate and release is closer to the young-cell limit $p\to\delta$; later, if
the epidemic grows, generation-time structure (and hence the operative
$p_{\rm eff}(r)$) departs from the $r=0$ value. The relative discrepancy
\begin{equation}
\frac{\bigl|\Vfree^{\rm NEW}(t)-\Vfree^{\rm CLASSICAL}(t)\bigr|}
{\max_s \Vfree^{\rm NEW}(s)}
\end{equation}
is plotted in Figure~\ref{fig:rel-diff}. In supercritical regimes the error is
often order-one or larger by the end of the plotted window; in the high-death
panel it is especially severe because $V_\infty$ is small while the
age-structure of release remains strongly non-exponential.

\begin{figure}[H]
\centering
\includegraphics[width=\textwidth]{figures/overlay_rel_diff.pdf}
\caption{Relative free-virion discrepancy
$|\text{NEW}-\text{CLASSICAL}|/\max(\text{NEW})$.
Large values mean the classical constant-$p$ ODEs fail to track the new model.}
\label{fig:rel-diff}
\end{figure}

\begin{figure}[H]
\centering
\includegraphics[width=\textwidth]{figures/overlay_growth_phase.pdf}
\caption{Growth-phase mismatch for one supercritical set
($\beta=1$, $\mu=0$, $\delta=0.1$).
\textbf{Solid navy = NEW model};
\textbf{dashed orange = CLASSICAL} with $r{=}0$ matching;
\textbf{dash-dot purple = CLASSICAL} with young-cell limit
$p=\delta$, $d_{\Icell}=\mu+\delta$.
Neither classical choice tracks the new curve over the whole trajectory.
Left panel logarithmic in $\Vfree$.}
\label{fig:growth-phase}
\end{figure}

\paragraph{Reading the plots.}
\begin{itemize}
\item Always: \textbf{solid navy = NEW}, \textbf{dashed orange = CLASSICAL}.
\item \textbf{CLASSICAL matched} ($p=p_{\rm eff}(0)$) is the fairest
constant-parameter comparison. It often gets the order of magnitude, but not
the shape (early lag, peak timing, late growth rate).
\item \textbf{CLASSICAL young-cell} ($p=\delta$) is better only briefly, then
under-produces.
\item \textbf{CLASSICAL naive} ($p=\langle X\rangle_{\rm QS}$) treats a count as a
rate and is systematically wrong in scale
(Figure~\ref{fig:overlay-V-naive}).
\item Discrepancies grow with supercritical $R_0$ and with strong
age-structure in the kernels (slow bursting, nonzero $\mu$).
\end{itemize}
These numerical facts are exactly what the analytic reduction of
Section~\ref{sec:peff} predicts: constant $(p,d_{\Icell})$ can match one growth
rate~$r$ exactly, but not an entire transient that samples a continuum of
ages.

% ===========================================================================
\section{Effective parameters, $R_0$, and growth-phase dependence}
\label{sec:peff}
% ===========================================================================

\subsection{Mean productive lifetime}

\begin{result}
\begin{equation}
\E[T_{\rm prod}]
=\int_0^\infty\Ihat(t)\,\dd t
=\frac{1}{\beta}\log\!\Bigl(\frac{a}{a-1}\Bigr)
=\frac{1}{\beta}\log\langle X\rangle_{\rm QS}.
\end{equation}
\end{result}
When $\mu=0$ this is $\beta^{-1}\log(1+\beta/\delta)$. The productive lifetime
is the time exponential growth at rate~$\beta$ takes to reach the
quasi-stationary level.

\subsection{Exact reduction in the exponential phase}

Substitute $i(t)=i_0\ee^{rt}$ into the renewal system and compare with the
classical BMVR. Writing $\Lap{f}(r)=\int_0^\infty\ee^{-r\alpha}f(\alpha)\,\dd\alpha$:
\begin{result}[Effective BMVR parameters]
\begin{equation}
p_{\rm eff}(r)=\frac{\delta\,\Lap{K}(r)}{\Lap{\Ihat}(r)},
\qquad
d_{\Icell,{\rm eff}}(r)=\frac{1}{\Lap{\Ihat}(r)}-r.
\end{equation}
\end{result}
This reduction is \emph{exact} for exponential solutions (not for transients):
substituting into the BMVR eigenvalue relation recovers the renewal
characteristic equation. Limits:
\begin{align}
r=0:&\quad
p_{\rm eff}=\frac{\beta\,V_\infty}{\log\bigl(a/(a-1)\bigr)},
\quad
d_{\Icell,{\rm eff}}=\frac{\beta}{\log\bigl(a/(a-1)\bigr)};\\
r\to\infty:&\quad
p_{\rm eff}\to\delta,\qquad
d_{\Icell,{\rm eff}}\to\mu+\delta.
\end{align}
The young-cell limit is intuitive: a newly infected cell holds one particle and
bursts at rate~$\delta$.

\subsection{\texorpdfstring{$R_0$}{R0} is unchanged by bursting}

\begin{result}
The characteristic equation is
\[
r+c=\gamma T\,\delta\Lap{K}(r),
\qquad\text{hence}\qquad
R_0=\frac{\gamma T\,V_\infty}{c}.
\]
At fixed total per-cell output $V_\infty$, replacing constant budding by bursting
leaves $R_0$ \emph{exactly} unchanged. Bursting alters generation times ---
hence the growth rate~$r$ and the extinction probability --- but not the
threshold.
\end{result}

\subsection{\texorpdfstring{$p$}{p} is not a constant of the pathogen}

Because $p_{\rm eff}$ decreases from $p_{\rm eff}(0)$ toward $\delta$ as $r$
grows, the effective release rate is smaller during acute exponential growth
(young cells, low loads) than at set-point. At
$(\beta,\mu,\delta)=(1,0,0.1)$ one has $p_{\rm eff}(0)\approx 4.59$ against
$p_{\rm eff}(\infty)=0.1$.

\begin{result}[Prediction]
BMVR fits to acute-phase data and to set-point data should yield systematically
different~$p$, in a direction and by an amount the theory specifies.
Equivalently: the eclipse-like delay usually inserted by hand into within-host
models emerges here from the intracellular dynamics, with no extra parameter.
\end{result}

\subsection{Identifiability (a useful negative result)}

Three intracellular rates $(\beta,\mu,\delta)$ map onto two BMVR parameters
$(p,d_{\Icell})$ at any given~$r$. The map is not invertible. Recovering the
intracellular rates therefore needs a third observable, for example:
\begin{enumerate}
\item burst-size dispersion (the geometric ratio $1/a$);
\item fraction of infected cells that clear without bursting ($b$);
\item shape of the release-time density $\varphi=\delta J$.
\end{enumerate}
Any one of these plus $(p,d_{\Icell})$ over-determines $(\beta,\mu,\delta)$ and
gives a consistency test.

% ===========================================================================
\section{The flooding advantage}
% ===========================================================================

This is the strongest single biological result: it answers, exactly, when
bursting helps a pathogen establish at low numbers.

\subsection{Offspring of one infected cell}

Each released particle independently infects with probability
$q=\gamma T/(\gamma T+c)$. The offspring generating function of one bursting
cell is
\begin{equation}
G_{\rm off}(z)=b+\frac{\delta}{\beta}\,\frac{y}{a-y},
\qquad y=1-q+qz,
\end{equation}
with mean $m=qV_\infty$ ($=R_0$ when $\gamma T\ll c$).

\subsection{Extinction probability in closed form}

$G_{\rm off}(z)=z$ is quadratic in~$y$. The non-trivial root is
\begin{equation}
z_{\rm ext}^{\rm burst}
=\frac{a(1-q+qb)-1+q}{q}.
\end{equation}

\subsection{\texorpdfstring{Comparison with matched-$R_0$ budding}{Comparison with matched-R0 budding}}

Comparator: a BMVR cell with exponential lifetime producing at constant rate,
matched so that total mean output equals $V_\infty$, with the same~$q$. Its
offspring number is geometric with mean~$R_0$, so
$z_{\rm ext}^{\rm bud}=1/R_0$.

Write $L:=a(1-b)=a-\mu/\beta$, so $V_\infty=L/(a-1)$. Then:
\begin{result}[Flooding advantage --- main theorem]
\label{res:flood}
For $R_0>1$,
\begin{equation}
z_{\rm ext}^{\rm burst}-z_{\rm ext}^{\rm bud}
=\bigl(L-1\bigr)\Bigl(\frac{1}{R_0}-1\Bigr).
\end{equation}
Hence bursting strictly lowers extinction probability relative to matched-$R_0$
budding if and only if $L>1$, i.e.
\begin{equation}
\boxed{\;
\frac{1}{\delta}\;<\;\frac{1}{\beta}+\frac{1}{\mu}\;}.
\end{equation}
When $L=1$ the two extinction probabilities are \emph{identical for every~$q$}.
When $\mu=0$ the condition holds for all $\delta>0$: bursting always wins.
\end{result}

\paragraph{Interpretation.}
Bursting concentrates output into one batch, which lowers offspring variance
relative to the geometric output of a budding cell. But it also carries the
all-or-nothing hazard that the intracellular population dies out first
(probability~$b$) and the cell releases nothing. The condition says the
advantage survives exactly when the burst timescale $1/\delta$ is shorter than
the harmonic-type combination of birth and death timescales.

\paragraph{Numerical checks (branching simulation).}
\begin{center}
\renewcommand{\arraystretch}{1.2}
\begin{tabular}{@{}llrrr@{}}
\toprule
$(\beta,\mu,\delta)$ & $q$ & $z^{\rm burst}$ & $z^{\rm bud}$ & Winner \\
\midrule
$(1,0,0.1)$ & $0.2$ & $0.400$ & $0.455$ & bursting ($L>1$) \\
$(1,0.9,0.1)$ & $0.9$ & $0.935$ & $0.844$ & budding ($L<1$) \\
$(1,0.5,1/3)$ & $0.6$ & $0.833$ & $0.833$ & tie ($L=1$) \\
\bottomrule
\end{tabular}
\end{center}

\begin{remark}
The comparison holds the offspring \emph{mean} fixed and is therefore a
statement about branching structure only. It says nothing about nonlinear
effects (immune saturation, local depletion) of the kind Komarova's ``all run
out at once'' intuition concerns. Both mechanisms should be named and
separated in any write-up.
\end{remark}

% ===========================================================================
\section{Further consequences}
% ===========================================================================

\subsection{Multiplicity of infection}

For a cell founded by $k$ particles,
\[
\Ihat_k(a)=I(a)^k-D(a)^k,\qquad
g_k(a)=\delta\bigl[k\,K(a)+k(k-1)J(a)^2\bigr].
\]
The release kernel is \emph{superlinear} in~$k$, so MOI matters even in the
mean --- a prediction the classical BMVR cannot express, and one that single-cell
or low-MOI assays can test.

\subsection{Two-type process and the eclipse phase}

Chapter~5 already gives the exact finite-time no-catastrophe probability $S(t)$
for a two-type BDC with type-specific rates $(\delta_1,\delta_2)$. The regime
$\delta_1=0<\delta_2$ (GATE: no burst before conversion) has $S'(0)=0$ --- a
zero initial hazard followed by a rising one. That is precisely the eclipse-phase
signature that within-host models normally insert by hand; here it emerges
mechanistically, with eclipse duration set by the conversion rate. Completing
the two-type release and survival kernels is the natural next technical step
and would make Chapters~4 and~5 one continuous story.

\subsection{Single-cell assay predictions (lytic BDC)}

A limiting-dilution assay of the Hataye type measures single-lineage release.
For a pure one-type BDC, predictions include:
\begin{enumerate}
\item a fraction $b$ of wells show no release ever;
\item cumulative release per well is a step: zero, then one terminal jump;
\item jump times follow density $\delta J(t)/(1-b)$;
\item jump sizes are geometric with ratio $1/a$ (same law as the QSD);
\item between-well variance equals $\mathrm{Var}(W_\infty)$, which the
deterministic BMVR cannot produce.
\end{enumerate}
Those predictions are appropriate for genuinely lytic systems
(\textit{Y.~pestis}, \textit{F.~tularensis}, lytic phage). For HIV they are
too strong: see the next section.

% ===========================================================================
\section{Distinction with HIV: what Hataye et al.\ show, and a candidate model}
\label{sec:hiv}
% ===========================================================================

The core BDC and renewal results above are written for pathogens whose
intracellular life ends in a \textbf{load-proportional catastrophe} (lyse and dump).
HIV is different. It \textbf{buds} from the plasma membrane; the infected cell is not
required to rupture and empty itself in one event. Single-cell and limiting-dilution
data nevertheless show something that \emph{looks} surge-like: long quiet periods,
then a concentrated bout of release. The right modelling response is not to force HIV
into pure BDC, nor to keep classical constant-$p$ BMVR, but to use an
\textbf{eclipse-structured, short productive pulse} with the same renewal free-virus
coupling.

\subsection{What the Hataye et al.\ (2019) paper shows}

Hataye, Casazza, Best, Perelson, Koup and co-authors
(\textit{Cell Host \& Microbe} 26:748--763, 2019) couple limiting-dilution
latency-reactivation cultures to stochastic models. The findings that matter for
equation choice are the following.

\paragraph{On initial release (viral inhibition cultures; no de novo infection).}
\begin{enumerate}
\item Timing and magnitude of first HIV RNA release from reactivated latently
infected CD4 T cells are \textbf{highly variable} (total yield, delay to first
detection, detection duration).
\item Mean release is delayed: first detection most often around day~3, with a
\textbf{sudden bulk of production around days 3--5} --- not a flat continuous drip
from day~0.
\item A simple productive-cell model (release at constant rate $p$, death at rate
$\delta$, total yield geometric with mean $p/\delta$) is
\textbf{under-dispersed} relative to the data (Fano factor too small).
\item A single exponential eclipse $E\to I$ fits the delay poorly (mode at day~1).
A \textbf{multi-stage eclipse} (Erlang-like chain of $n$ compartments; they favour
$n\approx 5$) fixes the mean timing.
\item Multi-stage eclipse alone still under-disperses. Transcriptional
ON/OFF toggling does not fix dispersion under realistic CD4 division/death rates
($\rho\approx\mu\approx 0.5/$day from CFSE and population data).
\item Their best structure (\textbf{Model 3}) uses \textbf{two distinct eclipse
paths} (high- and low-release potential), multi-stage delay, and division/death
in eclipse. Then:
\begin{itemize}
\item many stimulated lineages release \textbf{nothing} detectable
  ($f_{\mathrm{HIV}}\approx 0.41$ of single LICs give detectable virus in their fit);
\item most productive sojourns in $I$ last on the order of \textbf{about a day};
\item multi-day ``sustained'' detection is largely from \textbf{proliferation in
  eclipse} then staggered short $E\to I$ transitions of descendants --- not one
  cell budding steadily for a week.
\end{itemize}
\end{enumerate}

\paragraph{On establishment (outgrowth cultures; de novo infection allowed).}
\begin{enumerate}
\item Release of replication-competent virus often fails to establish exponential
spread.
\item Establishment probability as a function of the expected number of
virus-releasing LICs per well is \textbf{sigmoid}, not the concave curve of
independent per-cell success. They interpret this as an \textbf{Allee effect}
(synergy at low census).
\item A critical threshold sits near $\Lambda\approx 2.3$ LICs, corresponding on
average to roughly \textbf{$5100$ detected HIV RNA copies} of initial release
(with their establishment definition $\ge 2\times 10^5$ RNA copies).
\item For exactly one stimulated LIC, their fitted dynamics predict establishment
only on the order of \textbf{$\sim 2\%$} of the time --- consistent with the known
gap between intact provirus and quantitative viral outgrowth.
\end{enumerate}

\paragraph{Reading for this document.}
Hataye et al.\ therefore support three modelling decisions and reject two others:
\begin{center}
\renewcommand{\arraystretch}{1.3}
\begin{tabular}{@{}p{0.46\textwidth}p{0.46\textwidth}@{}}
\toprule
Supported for HIV & Not supported as literal HIV \\
\midrule
Multi-stage eclipse (delay) & Pure one-type BDC terminal catastrophe \\
Short productive pulse then cell death & Memoryless constant-$p$ BMVR from $t=0$ \\
High stochasticity; many silent lineages & Independent branching alone near threshold \\
Renewal / stage structure for free virus & QSD $=$ geometric burst size as an HIV law \\
Allee / synergistic establishment & \\
\bottomrule
\end{tabular}
\end{center}
The visual ``surge'' is real, but in their best model it is
\textbf{delay + brief constant-$p$ production + latent lineage dynamics}, not a
single lytic dump of intracellular load.

\subsection{Biological distinction from the lytic BDC}

\begin{center}
\renewcommand{\arraystretch}{1.25}
\begin{tabular}{@{}lll@{}}
\toprule
 & Lytic BDC (this document, core) & HIV (Hataye-guided) \\
\midrule
Release mechanism & Catastrophe at rate $\delta X$ & Budding while productive \\
End of production & Terminal; process absorbed & Short $I$ lifetime; cell dies \\
Intracellular load $X$ & Explicit birth--death state & Not required as primary state \\
Quiet period & Until first catastrophe & Multi-stage eclipse (+ latency) \\
Sustained multi-day RNA & One large dump (or nothing) & Often staggered short pulses \\
Establishment theory & Independent offspring OK as baseline & Allee synergy near threshold \\
Natural hosts & \textit{Y.\ pestis}, \textit{F.\ tularensis}, phage & CD4 T cell latency / rebound \\
\bottomrule
\end{tabular}
\end{center}

\subsection{Candidate equations for the HIV case}
\label{sec:hiv-eqns}

The following is a compact, Hataye-aligned system suitable for coupling to the
renewal free-virus layer of Result~\ref{res:new-vde}. Notation is chosen to
avoid collisions with the BDC rates $(\beta,\mu,\delta)$.

\subsubsection*{Latent / eclipse / productive stages (one LIC lineage)}

Let $L$ be a stimulated latently infected cell (or the start of reactivation).
Two eclipse pathways $A$ (high release potential) and $B$ (low, brief) run in
parallel; pathway $A$ is multi-stage.

\begin{empheq}[box=\fbox]{align}
L &\;\xrightarrow{\alpha_A}\; E_1^{(A)},
&
L &\;\xrightarrow{\alpha_B}\; E^{(B)},
\label{eq:hiv-split}\\[0.35em]
E_j^{(A)}
&\;\xrightarrow{a}\; E_{j+1}^{(A)},
\quad j=1,\dots,n-1,
&
E_n^{(A)}
&\;\xrightarrow{a}\; I^{(A)},
\label{eq:hiv-erlang}\\[0.35em]
E^{(B)}
&\;\xrightarrow{a_B}\; I^{(B)}.
&&
\label{eq:hiv-B}
\end{empheq}
In each eclipse compartment, cells may divide or die without releasing virus:
\begin{equation}
\label{eq:hiv-eclipse-bd}
E \;\xrightarrow{\rho}\; E+E,
\qquad
E \;\xrightarrow{\mu_E}\; \varnothing
\qquad
\text{(same rates on $E_j^{(A)}$ and, if desired, on $E^{(B)}$)}.
\end{equation}
Hataye et al.\ take $\rho\approx\mu_E\approx 0.5\,/\mathrm{day}$ from CFSE and
bulk population stability under stimulation. Productive cells release free virus
(or supernatant RNA) at a constant rate and die:
\begin{empheq}[box=\fbox]{align}
I^{(A)} &\;\xrightarrow{p_A}\; I^{(A)}+\Vfree,
&
I^{(A)} &\;\xrightarrow{\delta_I}\; \varnothing,
\label{eq:hiv-prodA}\\[0.25em]
I^{(B)} &\;\xrightarrow{p_B}\; I^{(B)}+\Vfree,
&
I^{(B)} &\;\xrightarrow{\delta_I}\; \varnothing,
\label{eq:hiv-prodB}
\end{empheq}
with $p_A\gg p_B$ (or $p_B$ small and $I^{(B)}$ short-lived) to match the
bimodal release distribution. Free virus / RNA in culture or plasma decays
\begin{equation}
\label{eq:hiv-clear}
\Vfree \;\xrightarrow{c}\; \varnothing.
\end{equation}
(In their inhibition cultures $c$ is set by measured supernatant RNA half-life;
in vivo $c$ is the usual plasma clearance.)

\paragraph{Mean-field stage ODEs (large $\Lambda$).}
Writing $E_j$ for pathway $A$ and suppressing the dual path for display,
\begin{align}
\ddt{E_1}
&= \alpha_A L_0\,\delta_0(t)
+\rho E_1
-(\mu_E+\rho+a)E_1,
\label{eq:hiv-ode-E1}\\
\ddt{E_j}
&= a E_{j-1}
+\rho E_j
-(\mu_E+\rho+a)E_j,
\quad j=2,\dots,n,
\label{eq:hiv-ode-Ej}\\
\ddt{I}
&= a E_n
-\delta_I I,
\label{eq:hiv-ode-I}\\
\ddt{\Vfree}
&= p_A I
-c\,\Vfree
\qquad\text{(plus the $B$-path contribution $p_B I^{(B)}$)}.
\label{eq:hiv-ode-V}
\end{align}
(The Dirac source stands for an initial cohort; for ongoing infection replace it
by an incidence term as below.) Stochastic simulation of the same transitions
(Gillespie) is what Hataye et al.\ use to match single-well dispersion.

\subsubsection*{Coupling to free virus and new infections}

When de novo infection is allowed, incidence should feed new eclipse cells, not
instant productively infected cells:
\begin{empheq}[box=\fbox]{align}
i(t)
&= f\bigl(\Vfree(t),T\bigr),
\label{eq:hiv-incidence}\\[0.25em]
\text{source into }E_1^{(A)}
&= \pi_A\,i(t),
\qquad
\text{source into }E^{(B)}
&= \pi_B\,i(t),
\label{eq:hiv-split-inc}
\end{empheq}
with $\pi_A+\pi_B=1$ the pathway split.

\paragraph{Linear infection (classical).}
\[
f(\Vfree,T)=\gamma T\,\Vfree.
\]
This is the right large-population limit and recovers a stage-structured BMVR.

\paragraph{Allee / synergistic infection (Hataye establishment).}
Near the critical threshold their data favour a sigmoid establishment law, not
independent per-virion success. A minimal phenomenological upgrade is
\begin{equation}
\label{eq:hiv-allee}
f(\Vfree,T)
=\gamma T\,\Vfree\cdot
\frac{\Vfree}{K_A+\Vfree},
\end{equation}
so that infection is inefficient at very low $\Vfree$ and saturates to the usual
mass-action form when $\Vfree\gg K_A$. The scale $K_A$ is chosen so that
cumulative initial release near $\sim 5\times 10^3$ RNA copies sits at the
inflection of establishment probability (Hataye's empirical threshold). More
mechanistic forms (cooperative entry, restriction-factor saturation, local
multiplicity) can replace \eqref{eq:hiv-allee}; the point is that
\textbf{linearity in $\Vfree$ alone is not faithful near establishment}.

\subsubsection*{Renewal kernels induced by the HIV stage model}

Exactly as in the lytic case, the population free-virus equation is a
convolution against an age kernel --- only the kernel changes. Let $g_{\mathrm{HIV}}(a)$
be the expected release rate at age $a$ after a single infection (or reactivation),
and $S_{\mathrm{HIV}}(a)$ the probability of still being in a productive (or
not-yet-cleared) lineage. Then the same structural system as
Result~\ref{res:new-vde} holds:
\begin{empheq}[box=\fbox]{align}
\Icell(t)
&= \Icell_0 S_{\mathrm{HIV}}(t)
+\bigl(i * S_{\mathrm{HIV}}\bigr)(t),
\label{eq:hiv-Icell}\\[0.35em]
\ddt{\Vfree}
&= \Icell_0\,g_{\mathrm{HIV}}(t)
+\bigl(i * g_{\mathrm{HIV}}\bigr)(t)
-c\,\Vfree(t).
\label{eq:hiv-Vfree}
\end{empheq}
For the stage model above, $g_{\mathrm{HIV}}$ is \emph{not} $\delta K$ from the BDC.
It is the density of productive sojourn times times $p_A$ (and $p_B$), with the
Erlang eclipse supplying the delay. In particular $g_{\mathrm{HIV}}(a)\approx 0$ for
small $a$, then rises and falls --- the mathematical expression of
``nothing for a long time, then a lot.''

The effective parameters of Section~\ref{sec:peff} carry over with the
substitution $(\Ihat,g)\mapsto(S_{\mathrm{HIV}},g_{\mathrm{HIV}})$:
\begin{equation}
p_{\rm eff}(r)
=\frac{\Lap{g_{\mathrm{HIV}}}(r)}{\Lap{S_{\mathrm{HIV}}}(r)},
\qquad
R_0
=\frac{\gamma T}{c}\int_0^\infty g_{\mathrm{HIV}}(a)\,\dd a
\quad\text{(linear infection)}.
\end{equation}
Under \eqref{eq:hiv-allee}, $R_0$ remains a large-population threshold; early
establishment is controlled by the Allee scale $K_A$ as well.

\subsection{What to keep from the BDC theory, and what not to claim for HIV}

\begin{itemize}
\item \textbf{Keep:} renewal free-virus coupling; age/stage kernels; the idea that
constant $p$ is only an exponential-phase projection; stochastic early failure;
eclipse as a real dynamical object (GATE / multi-stage).
\item \textbf{Do not claim for HIV:} terminal catastrophe at rate $\delta X$;
QSD $=$ geometric burst size; pure one-shot cumulative release as the only
single-cell trajectory; independent branching as the full establishment story
near the Hataye threshold.
\item \textbf{Flooding comparison:} still meaningful as a \emph{linear} baseline
(matched mean yield, compare variance of offspring). Hataye adds that near
$\sim 5\times 10^3$ RNA copies, \textbf{synergy can dominate} variance effects.
Any HIV write-up should separate the two mechanisms.
\end{itemize}

\subsection{Suggested role in a paper or thesis chapter}

\begin{enumerate}
\item Core publishable BDC + renewal BMVR results (lytic end of the spectrum).
\item One section or companion note: \textbf{HIV is not BDC}; Hataye Model~3
as the empirical intermediate; equations
\eqref{eq:hiv-split}--\eqref{eq:hiv-Vfree}.
\item Numerically extract $g_{\mathrm{HIV}}(a)$ from the stage model (or from
Hataye's daily release curves) and overlay against constant-$p$ BMVR exactly as
in Section~\ref{sec:overlays}.
\item Establishment: independent branching vs Allee form \eqref{eq:hiv-allee},
with the critical release scale taken from Hataye et al.
\end{enumerate}
That keeps the mathematical core honest for lytic pathogens while giving HIV a
faithful, equation-level treatment grounded in the paper's data and models.

% ===========================================================================
\section{Reset-catastrophe process and its renewal BMVR}
\label{sec:reset-bmvr}
% ===========================================================================

This section is self-contained. It records an intermediate intracellular model
between pure continuous budding and absorbing BDC, and shows that it induces
\textbf{the same renewal BMVR structure} as Section~\ref{sec:overlays} / 
Result~\ref{res:new-vde}, with different kernels only.

\paragraph{Biological status (one sentence).}
Interpret $X$ as a local assembled pool (not whole-cell contents): full dump is a
mathematical extreme of batch release; HIV budding is better thought of as many
small exports, with this model as a tractable multi-pulse caricature on the
budding--bursting spectrum.

\subsection{Intracellular process (one productively infected cell)}

Age $a=0$ at the start of productive infection. Load $X_a\in\{0,1,2,\ldots\}$.
While the cell is alive:
\begin{center}
\renewcommand{\arraystretch}{1.2}
\begin{tabular}{@{}lll@{}}
\toprule
Event & Jump & Rate \\
\midrule
Birth & $n\to n+1$ & $\beta n$ \\
Death / degradation & $n\to n-1$ & $\mu n$ \\
Immigration (optional) & $n\to n+1$ & $\nu$ \\
\textbf{Dump (reset catastrophe)} & $n\to 0$, release $+n$ extracellular & $\delta n$ \\
\bottomrule
\end{tabular}
\end{center}
The cell itself dies at a separate rate $d_{\Icell}$ (cytopathic / immune),
independent of $X$ in the simplest case. On cell death the process stops.
Unlike absorbing BDC, a dump does \textbf{not} end the lineage: the pool
resets to $0$ and may re-accumulate, so one cell can release \textbf{many}
times before death.

\subsection{Kernels}

Dump intensity in state $n$ is $\delta n$ and the dump size is $n$, so the
expected release rate at age $a$ is $\delta\E[X_a^2\mathbf 1_{\{\text{alive}\}}]$.
With independent cell death at rate $d_{\Icell}$,
\begin{empheq}[box=\fbox]{equation}
\label{eq:g-reset}
S(a)=e^{-d_{\Icell}a},
\qquad
g(a)=\delta\,e^{-d_{\Icell}a}\,\E[X_a^2],
\end{empheq}
where $\E[X_a^2]$ is computed for the reset process (no absorption at a rupture
state). The structural formula $g\propto\E[X^2]$ matches BDC; the second moment
differs because $X$ restarts after each dump instead of ending in $R$.

Linear rates keep moments tractable: $\E[X_a]$ and $\E[X_a^2]$ obey a closed
ODE system under the process generator (birth, death, immigration, and the
reset jump). Then $g$ is explicit once those ODEs are solved.

\subsection{Renewal BMVR (same skeleton as BDC)}

With target cells $T$ fixed, infection rate $\gamma$, clearance $c$, and
incidence $i(t)=\gamma T\,\Vfree(t)$:
\begin{empheq}[box=\fbox]{align}
i(t)
&= \gamma T\,\Vfree(t),
\label{eq:reset-i}\\[0.35em]
\Icell(t)
&= \Icell_0\,S(t)
+\int_0^t i(t-\alpha)\,S(\alpha)\,\dd\alpha,
\label{eq:reset-I}\\[0.35em]
\ddt{\Vfree}
&= \Icell_0\,g(t)
+\int_0^t i(t-\alpha)\,g(\alpha)\,\dd\alpha
-c\,\Vfree(t),
\label{eq:reset-V}
\end{empheq}
with kernels \eqref{eq:g-reset}. In convolution form:
$\Icell=\Icell_0 S+i*S$ and $\Vfree'=\Icell_0 g+i*g-c\Vfree$.

This is identical in form to Result~\ref{res:new-vde}. Only $(S,g)$ change.

\subsection{Effective parameters and mature-cell limit}

For exponential incidence $i\propto e^{rt}$,
\begin{equation}
\label{eq:reset-peff}
p_{\rm eff}(r)=\frac{\Lap{g}(r)}{\Lap{S}(r)},
\qquad
d_{\Icell,{\rm eff}}(r)=\frac{1}{\Lap{S}(r)}-r.
\end{equation}
With $S(a)=e^{-d_{\Icell}a}$ one has $\Lap{S}(r)=1/(r+d_{\Icell})$, hence
$d_{\Icell,{\rm eff}}\equiv d_{\Icell}$ and
$p_{\rm eff}(r)=(r+d_{\Icell})\Lap{g}(r)$. The linear-infection threshold is
\begin{equation}
R_0=\frac{\gamma T}{c}\int_0^\infty g(a)\,\dd a
=\frac{\gamma T\cdot\E[\text{lifetime output per cell}]}{c}.
\end{equation}
Lifetime yield is finite only if $d_{\Icell}>0$ (otherwise recurrent dumps can
continue forever).

\paragraph{Mature-cell limit.}
If $X$ admits a stationary law $\pi$ on $\{0,1,\ldots\}$ and $d_{\Icell}$ is
small,
\[
p_\infty
=\lim_{a\to\infty}\frac{g(a)}{S(a)}
=\delta\,\E_\pi[X^2].
\]
A long-lived infected cell then behaves as classical constant-$p$ BMVR with
that $p$. Transients (first dumps, young cells) still require the full kernel
$g(a)$.

\subsection{Comparison with absorbing BDC}

\begin{center}
\renewcommand{\arraystretch}{1.25}
\begin{tabular}{@{}lll@{}}
\toprule
 & Absorbing BDC & Reset catastrophe \\
\midrule
After dump & process ends in $R$ & $X\to 0$, continues \\
Dumps per cell & $\le 1$ & many (until cell death) \\
Survival kernel & $\Ihat$ (no burst / not extinct) & typically $e^{-d_{\Icell}a}$ \\
Release kernel & $\delta K_{\mathrm{BDC}}(a)$ & $\delta e^{-d_{\Icell}a}\E[X_a^2]$ \\
Population equations & renewal BMVR & \textbf{same form} \\
Mature limit & not constant budding & $p\to\delta\E_\pi[X^2]$ \\
\bottomrule
\end{tabular}
\end{center}

\paragraph{Optional eclipse.}
Set $g(a)=0$ for $a<\tau_e$ (or Erlang-delay the start of the reset process);
keep $S(a)=e^{-d_{\Icell}a}$. Equations \eqref{eq:reset-i}--\eqref{eq:reset-V}
are unchanged.

\paragraph{Takeaway.}
Any intracellular model that supplies $(S,g)$ induces renewal BMVR. The
reset-catastrophe process does so with a multi-pulse release kernel and a
classical budding limit for mature cells, while remaining in the same
linear-rate / moment-closed analytic class as BDC.

% ===========================================================================
\section{HIV intracellular mean growth: linear vs exponential}
\label{sec:hiv-mean-growth}
% ===========================================================================

This section records a modelling decision for coarse-grained HIV intracellular
dynamics: whether mean load grows \emph{exponentially} (birth--death) or
\emph{linearly} (template-limited production), and what that implies for a
minimal regime. Export / extracellular release is left open.

\subsection{The dichotomy}

Let $m(a)=\E[X_a]$ be mean intracellular load (Gag, genomes, or
assembly-ready units) after productive infection begins.

\begin{center}
\renewcommand{\arraystretch}{1.25}
\begin{tabular}{@{}lll@{}}
\toprule
Regime & Mean law (schematic) & Mechanism \\
\midrule
\textbf{Mean-exponential} & $m'=r\,m$ & each unit generates more units (branching / BDC) \\
\textbf{Mean-linear} & $m'=\nu-\mu m$ (or $m'\approx\nu$) & production set by fixed templates / machinery \\
\bottomrule
\end{tabular}
\end{center}

Absorbing BDC and pure birth--death are exponential-mean. Immigration /
constant production (and the productive phase of classical BMVR) are
linear-mean.

\subsection{What HIV biology favours}

HIV does \textbf{not} amplify infectious units inside the cell by branching:

\begin{itemize}
\item After integration, the DNA template count is typically \textbf{one or a few
proviruses} --- not an exponentially growing genome pool.
\item Late genomic RNA is made by host Pol~II from that fixed DNA, not by
cytoplasmic RdRp-style genome cascading.
\item Gag and other structural products are translated from mRNA; they do not
self-replicate. Production is immigration-like: rate set by mRNA and ribosomes,
minus degradation / packaging / export.
\end{itemize}

So pure $m'=rm$ birth--death is a \textbf{poor} literal description of HIV
intracellular load. It fits phage genome races, intracellular bacterial
division, or some cytoplasmic RNA viruses better.

The main nonlinearity is \textbf{Tat positive feedback}: transcription rate
rises from a low basal level toward a saturated high state. During the switch,
product counts can rise faster than a pure linear ramp (even looking briefly
super-linear) because the \emph{production rate} is increasing --- not because
each unit is a branching particle. Once fully ON, production is well
approximated as roughly \textbf{constant} (or fluctuating about a high mean),
hence mean accumulation is linear / saturating.

\paragraph{Confidence.}
\begin{itemize}
\item \textbf{High:} not pure exponential-mean birth--death of virion-like units.
\item \textbf{Medium--high:} main productive window is closer to mean-linear
(template-limited) than exponential.
\item \textbf{Medium:} exact shape of $m(a)$ (linear vs saturating vs sigmoidal)
depends on species counted and on the Tat phase; few studies cleanly fit
$e^{rt}$ vs $t$ to single-cell intracellular load.
\end{itemize}

\subsection{Does the Tat intermediate exceed the late linear rate?}

Tat-controlled transcription is saturating in Tat:
\[
\text{rate}
= r_0 + r_{\max}\,f(\mathrm{Tat}),
\qquad f\text{ increasing, }f\to 1.
\]
The late fully-ON rate is the \textbf{ceiling}. The intermediate phase
\textbf{climbs toward} that ceiling; in the mean, it does \textbf{not}
systematically exceed it.

Caveats (different observables):
\begin{itemize}
\item Stochastic transcriptional bursts can spike above the \emph{mean} late rate
(noise, not a higher intermediate phase).
\item \textbf{Net} accumulation $\dd X/\dd t$ can peak mid-course if export turns
on late, even when production is still rising or flat --- that is export
balance, not Tat overshoot.
\end{itemize}
For coarse modelling, treat Tat as acceleration \textbf{up to} the late rate, not
a regime hotter than late ON.

\subsection{Recommended coarse-grained regime}

A fair, defensible minimal regime for HIV intracellular dynamics is:
\begin{empheq}[box=\fbox]{equation}
\textbf{eclipse}
\;\longrightarrow\;
\textbf{linear (immigration-type) accumulation of }X.
\end{empheq}
\begin{enumerate}
\item \textbf{Eclipse:} negligible late product / no infectious export. Prefer a
fixed or multi-step (Erlang) delay so first activity is not peaked at $t=0^+$.
\item \textbf{Productive phase:} $m'\approx\nu-\mu m$ (or $m'\approx\nu$ while
accumulating), with $\nu$ the late fully-ON production rate. Tat switch-on may
be folded into a slightly longer eclipse plus an abrupt start of $\nu$.
\item \textbf{Export:} undecided at this stage (constant rate, linear in $X$,
batch, \ldots). The linear-vs-exponential choice for \emph{intracellular}
accumulation does not force the export law.
\end{enumerate}

This matches fixed provirus templates, Hataye-type delay then productive window,
and the mean-linear assessment above. It is \textbf{not} absorbing BDC, and it
does not claim every molecular species is literally linear --- only that one
coarse stock after eclipse is a fair reduction.

\paragraph{Model ranking (HIV intracellular mean growth).}
\begin{center}
\renewcommand{\arraystretch}{1.2}
\begin{tabular}{@{}ll@{}}
\toprule
Model & Fit as HIV intracellular growth \\
\midrule
Birth--death / absorbing BDC ($m'=rm$) & poor as literal mechanism \\
Immigration + export/dump ($m'\approx\nu-\cdots$) & much better caricature \\
Eclipse then constant production rate $\nu$ & best simple default \\
Tat ON/OFF + immigration & best simple mechanistic default \\
\bottomrule
\end{tabular}
\end{center}

\paragraph{Takeaway.}
Coarse-grained HIV: \textbf{eclipse, then mean-linear intracellular production};
export to be chosen later. Do not default to exponential-mean BDC load growth
inside the cell. Tat adds a switch up to the late rate, not a systematic
overshoot of it.

% ===========================================================================
\section{Building BMVR piece by piece: five cases}
\label{sec:bmvr-pieces}
% ===========================================================================

This section lays out a modular route from the classic Basic Model of Viral
Replication (BMVR) to models with eclipse and with an explicit intracellular
stock. Target-cell count $T$ is held constant (simplified BMVR). Free-virion
clearance is $c$; infection rate is $\gamma$. Notation for population counts:
$\Icell$ (productively infected cells), $\Vfree$ (free virions), and where
needed $E$ (eclipse cells) and $Q$ (total intracellular load summed over
infected cells).

\paragraph{Modular programme.}
\begin{enumerate}
\item \textbf{Case 1 --- Classic BMVR.} Linear production with
\emph{immediate} release at constant budding rate $p$ per cell. No
intracellular stock, no delay.
\item \textbf{Case 2 --- Eclipse + classic release.} Same as Case~1, but new
infections enter a non-producing eclipse class before becoming productive.
(Standard in the literature.)
\item \textbf{Case 3 --- Intracellular stock + continuous export.} Linear
accumulation into a compartment $Q$; continuous release from that compartment
into $\Vfree$ (no eclipse in the main write-up; eclipse stacks as in Case~2).
\item \textbf{Case 4 --- Intracellular stock + proportional full clear.} Same
accumulation as Case~3, but release is by dump events that clear the pool at a
load-proportional rate (reset catastrophe at the mean-field level). No eclipse.
\item \textbf{Case 5 --- Eclipse + stock + proportional full clear.} Case~4
with an eclipse class: incidence feeds $E$, only productively infected cells
accumulate $Q$ and dump.
\end{enumerate}
Cases 3--5 are the ``accumulation + structured release'' layer: the
\emph{export law} and optional eclipse distinguish them; intracellular
production is immigration-type (mean-linear), not birth--death.

% ---------------------------------------------------------------------------
\subsection{Case 1 --- Classic BMVR}
% ---------------------------------------------------------------------------

\paragraph{Assumptions.}
An infected cell produces virions that are released \emph{immediately} at constant
rate $p$ (budding). No eclipse, no stored intracellular pool. Infected cells are
removed at rate $d_{\Icell}$.

\begin{empheq}[box=\fbox]{align}
\ddt{\Icell}
&= \gamma T\,\Vfree - d_{\Icell}\,\Icell,
\label{eq:c1-I}\\[0.3em]
\ddt{\Vfree}
&= p\,\Icell - c\,\Vfree.
\label{eq:c1-V}
\end{empheq}

\paragraph{Reading the equations.}
Infection creates productively infected cells at rate $\gamma T\Vfree$. Each such
cell contributes $p$ virions per unit time to the free pool; free virions clear at
rate $c$. Production and export are collapsed into the single parameter $p$.

\paragraph{Basic reproduction number (linearised at $\Vfree$ small).}
\[
R_0=\frac{\gamma T\,p}{c\,d_{\Icell}}.
\]

% ---------------------------------------------------------------------------
\subsection{Case 2 --- Eclipse phase + classic release}
% ---------------------------------------------------------------------------

\paragraph{Assumptions.}
New infections enter an eclipse (or latent-productive delay) class $E$: infected
but not yet producing. Conversion $E\to\Icell$ at rate $a$ (mean eclipse length
$1/a$). Optional eclipse death $d_E$; productive cells as in Case~1. Release is
still immediate constant budding $p$ from $\Icell$ only.

This extension is standard in within-host HIV modelling (delay before
production; fits to acute viremia and treatment timing).

\begin{empheq}[box=\fbox]{align}
\ddt{E}
&= \gamma T\,\Vfree - (a+d_E)\,E,
\label{eq:c2-E}\\[0.3em]
\ddt{\Icell}
&= a\,E - d_{\Icell}\,\Icell,
\label{eq:c2-I}\\[0.3em]
\ddt{\Vfree}
&= p\,\Icell - c\,\Vfree.
\label{eq:c2-V}
\end{empheq}

\paragraph{Reading the equations.}
Incidence feeds $E$, not $\Icell$. Only after conversion do cells contribute to
$\Vfree$ at rate $p$. Setting $a\to\infty$ (instant exit from eclipse) and
$d_E=0$ recovers Case~1 with the same $p$.

\paragraph{$R_0$.}
A fraction $a/(a+d_E)$ of eclipse cells become productive; each then yields
mean $p/d_{\Icell}$ virions, so
\[
R_0
=\frac{\gamma T}{c}\cdot\frac{a}{a+d_E}\cdot\frac{p}{d_{\Icell}}.
\]

\paragraph{Multi-stage eclipse.}
Replace one $E$ by a chain $E_1\to\cdots\to E_n\to\Icell$ (Erlang delay) for a
more realistic distribution of eclipse times; the $\Vfree$ equation is unchanged.
Hataye-type multi-compartment eclipse is this idea at finer grain.

% ---------------------------------------------------------------------------
\subsection{Case 3 --- Intracellular accumulation + continuous export}
% ---------------------------------------------------------------------------

\paragraph{Assumptions.}
Productively infected cells accumulate an intracellular stock (total pool $Q$
across all such cells) by \textbf{linear production} at rate $\nu$ per cell
(immigration-type; Section~\ref{sec:hiv-mean-growth}). Virions are released
\textbf{continuously} from that stock: each intracellular unit is exported at
per-capita rate $\varepsilon$ (linear budding from the compartment). Cell death
removes both the cell and its share of load (mean-field: load lost at rate
$d_{\Icell}Q$). New infections start with empty stock.

\begin{empheq}[box=\fbox]{align}
\ddt{\Icell}
&= \gamma T\,\Vfree - d_{\Icell}\,\Icell,
\label{eq:c3-I}\\[0.3em]
\ddt{Q}
&= \nu\,\Icell - \varepsilon\,Q - d_{\Icell}\,Q,
\label{eq:c3-Q}\\[0.3em]
\ddt{\Vfree}
&= \varepsilon\,Q - c\,\Vfree.
\label{eq:c3-V}
\end{empheq}

\paragraph{Reading the equations.}
\begin{itemize}
\item $\Icell$: same bookkeeping as Case~1.
\item $Q$: gained by production $\nu\Icell$; lost by export $\varepsilon Q$ and
by death of infected cells $d_{\Icell}Q$.
\item $\Vfree$: fed only by export $\varepsilon Q$, not by a phenomenological
$p\Icell$.
\end{itemize}
Mean load per productive cell is $m=Q/\Icell$ (when $\Icell>0$). In a
quasi-steady intracellular balance ($\ddt Q\approx 0$),
\[
Q_*
=\frac{\nu}{\varepsilon+d_{\Icell}}\,\Icell,
\qquad
p_{\mathrm{eff}}
:=\frac{\varepsilon Q_*}{\Icell}
=\frac{\varepsilon\nu}{\varepsilon+d_{\Icell}},
\]
and Case~3 collapses to Case~1 with $p=p_{\mathrm{eff}}$. Thus classic BMVR is
the \textbf{fast-export / quasi-steady stock} limit of this model. When stock
dynamics are not fast, $\Vfree$ lags and depends on the history of $Q$.

\paragraph{Constant release rate per cell (variant).}
If instead each productive cell exports at a fixed rate $\rho$ independent of
load (only while $Q$ is not limiting), one may close with
$\ddt{\Vfree}=\rho\Icell-c\Vfree$ and drain $Q$ by $\rho\Icell$, provided
$Q$ stays non-negative --- a cruder law, closer to Case~1, with $Q$ tracking
unused stock. The linear law $\varepsilon Q$ is the natural default for
``release from the compartment.''

\paragraph{With eclipse (Case 2$+$Case 3).}
Incidence enters $E$; only $\Icell$ produces into $Q$:
\begin{align}
\ddt{E}
&= \gamma T\,\Vfree - (a+d_E)\,E,
\label{eq:c3e-E}\\
\ddt{\Icell}
&= a\,E - d_{\Icell}\,\Icell,
\label{eq:c3e-I}\\
\ddt{Q}
&= \nu\,\Icell - \varepsilon\,Q - d_{\Icell}\,Q,
\label{eq:c3e-Q}\\
\ddt{\Vfree}
&= \varepsilon\,Q - c\,\Vfree.
\label{eq:c3e-V}
\end{align}
Eclipse cells do not contribute to $Q$ or $\Vfree$ until conversion.

\paragraph{$R_0$ (Case 3 without eclipse).}
Lifetime output per cell equals production integrated against survival,
$\nu/d_{\Icell}$, if every produced unit is eventually either exported or lost
at death; more carefully, the export fraction is
$\varepsilon/(\varepsilon+d_{\Icell})$, so mean virions per cell
$=\nu\varepsilon/(d_{\Icell}(\varepsilon+d_{\Icell}))\cdot d_{\Icell}/\varepsilon$
simplifies to
\[
\frac{\nu}{d_{\Icell}}\cdot\frac{\varepsilon}{\varepsilon+d_{\Icell}},
\qquad
R_0
=\frac{\gamma T\,\nu\,\varepsilon}
{c\,d_{\Icell}\,(\varepsilon+d_{\Icell})},
\]
matching Case~1 with $p=\varepsilon\nu/(\varepsilon+d_{\Icell})$.

% ---------------------------------------------------------------------------
\subsection{Case 4 --- Intracellular accumulation + proportional full clear}
% ---------------------------------------------------------------------------

\paragraph{Assumptions.}
Same linear production into the stock as Case~3 ($\nu$ per productive cell).
Release is \textbf{not} continuous per unit. Instead, dumps occur at a rate
proportional to load: in a cell with load $x$, dump rate $\delta x$, and the
whole pool $x$ is transferred to $\Vfree$ while the cell's stock resets to $0$
(cell may survive; death is still the separate rate $d_{\Icell}$).

\paragraph{Mean-field closure.}
Let $Q$ be total intracellular load and $\Icell$ the number of productive cells.
Under a homogeneous-load closure (every cell carries $m=Q/\Icell$), the
population dump rate of virions is
\[
\delta\,\E\!\sum_i X_i^2
\approx
\delta\,\Icell\,m^2
=\delta\,\frac{Q^2}{\Icell},
\]
and the same quantity is the rate at which load leaves $Q$. (Exact stochastic
models use $\delta\E[X^2]$ per cell; this is the standard mean-field analogue.)

\begin{empheq}[box=\fbox]{align}
\ddt{\Icell}
&= \gamma T\,\Vfree - d_{\Icell}\,\Icell,
\label{eq:c4-I}\\[0.3em]
\ddt{Q}
&= \nu\,\Icell - \delta\,\frac{Q^2}{\Icell} - d_{\Icell}\,Q,
\label{eq:c4-Q}\\[0.3em]
\ddt{\Vfree}
&= \delta\,\frac{Q^2}{\Icell} - c\,\Vfree,
\label{eq:c4-V}
\end{empheq}
(with the convention that the nonlinear terms vanish when $\Icell=0$).

\paragraph{Reading the equations.}
\begin{itemize}
\item Production $\nu\Icell$ as in Case~3.
\item The term $\delta Q^2/\Icell$ is simultaneous \textbf{export into
$\Vfree$} and \textbf{clearing of $Q$} (full dump).
\item Cell death still removes load at rate $d_{\Icell}Q$; dumps do
\emph{not} remove the cell from $\Icell$ (reset, not lysis). For a lytic
variant, add a sink in $\Icell$ proportional to total dump rate.
\end{itemize}

\paragraph{Contrast with Case 3.}
Continuous export is linear in $Q$ ($\varepsilon Q$). Full clear is
\textbf{quadratic} in total load at fixed $\Icell$ ($\propto Q^2/\Icell$),
i.e.\ superlinear in mean load $m$. Mature quasi-steady balance
($\ddt Q\approx 0$, $\Icell$ slowly varying) gives
\[
\nu
=\delta m^2 + d_{\Icell}\,m
\quad\Rightarrow\quad
m
=\frac{-d_{\Icell}+\sqrt{d_{\Icell}^2+4\delta\nu}}{2\delta},
\]
and effective release per cell $p_{\mathrm{eff}}=\delta m^2$, which plays the
role of $p$ in Case~1 only in that steady regime --- transients and early
infection need the full $(Q,\Vfree)$ system. For the same release law
\textbf{with eclipse}, see Case~5.

\paragraph{$R_0$ (sketch).}
In the linearised regime near zero infection, $Q$ is small, dumps are rare, and
the first-order expansion of \eqref{eq:c4-Q}--\eqref{eq:c4-V} must be taken
carefully (the map $Q\mapsto Q^2/\Icell$ is singular at $0$ in naive form).
For invasion analysis it is safer to linearise the underlying single-cell
process (expected lifetime output $V_\infty$) and set
$R_0=\gamma T V_\infty/c$, as in the renewal framework. The ODE system
\eqref{eq:c4-I}--\eqref{eq:c4-V} is aimed at bulk dynamics with $\Icell$ not
near zero; use branching / renewal for establishment.

% ---------------------------------------------------------------------------
\subsection{Case 5 --- Eclipse + stock + proportional full clear}
% ---------------------------------------------------------------------------

\paragraph{Assumptions.}
Combine Case~2 and Case~4: new infections enter a non-producing eclipse class
$E$; only after conversion to productively infected $\Icell$ does linear
production into the intracellular stock $Q$ begin; release from $Q$ is by
load-proportional full clear (same mean-field dump as Case~4). Eclipse cells
carry no stock and do not feed $\Vfree$.

\begin{empheq}[box=\fbox]{align}
\ddt{E}
&= \gamma T\,\Vfree - (a+d_E)\,E,
\label{eq:c5-E}\\[0.3em]
\ddt{\Icell}
&= a\,E - d_{\Icell}\,\Icell,
\label{eq:c5-I}\\[0.3em]
\ddt{Q}
&= \nu\,\Icell - \delta\,\frac{Q^2}{\Icell} - d_{\Icell}\,Q,
\label{eq:c5-Q}\\[0.3em]
\ddt{\Vfree}
&= \delta\,\frac{Q^2}{\Icell} - c\,\Vfree,
\label{eq:c5-V}
\end{empheq}
(with nonlinear dump terms set to $0$ when $\Icell=0$).

\paragraph{Reading the equations.}
\begin{itemize}
\item \eqref{eq:c5-E}--\eqref{eq:c5-I}: same eclipse bookkeeping as Case~2.
  Incidence creates $E$; conversion at rate $a$ creates $\Icell$; optional
  eclipse death $d_E$; productive death $d_{\Icell}$.
\item \eqref{eq:c5-Q}--\eqref{eq:c5-V}: identical stock and release law to
  Case~4. Only productively infected cells produce ($\nu\Icell$) and dump
  ($\delta Q^2/\Icell$).
\item Setting $a\to\infty$ with $d_E=0$ recovers Case~4 (instant exit from
  eclipse). Setting the dump aside and replacing
  $\delta Q^2/\Icell$ by continuous export would give eclipse$+$Case~3.
\end{itemize}

\paragraph{What eclipse changes relative to Case 4.}
There is an additional delay before any contribution to $Q$ or $\Vfree$. Early
after a pulse of infection, $E$ rises first, then $\Icell$ and $Q$, then
$\Vfree$ --- a staged lag that Case~4 cannot produce. Quasi-steady relations
for mean load $m=Q/\Icell$ among \emph{productive} cells are unchanged from
Case~4 once $\Icell$ is slowly varying; the eclipse only filters how
$\Icell(t)$ is driven by $\Vfree$.

\paragraph{$R_0$ (sketch).}
As in Case~2, a fraction $a/(a+d_E)$ of eclipse cells become productive. If
$V_\infty$ is expected lifetime free-virus output of one productive cell under
the Case~4 release law, then
\[
R_0
=\frac{\gamma T}{c}\cdot\frac{a}{a+d_E}\cdot V_\infty
\]
in the linear-infection, single-cell sense. Bulk invasion near $\Icell=0$ has
the same caveat as Case~4 regarding the $Q^2/\Icell$ closure; prefer
single-cell / renewal linearisation for thresholds.

% ---------------------------------------------------------------------------
\subsection{Summary of the five cases}
% ---------------------------------------------------------------------------

\begin{center}
\renewcommand{\arraystretch}{1.3}
\begin{tabular}{@{}p{0.08\textwidth}p{0.38\textwidth}p{0.42\textwidth}@{}}
\toprule
Case & Intracellular story & Free-virus source term \\
\midrule
1 Classic & none (instant release) & $p\,\Icell$ \\
2 Eclipse + classic & delay $E\to\Icell$, then as 1 & $p\,\Icell$ \\
3 Stock + continuous export & $\nu$ into $Q$, export $\varepsilon Q$ & $\varepsilon\,Q$ \\
4 Stock + full clear & $\nu$ into $Q$, dump $\propto Q^2/\Icell$ & $\delta Q^2/\Icell$ \\
5 Eclipse + stock + full clear & $E\to\Icell$, then as 4 & $\delta Q^2/\Icell$ \\
\bottomrule
\end{tabular}
\end{center}

\paragraph{What is fixed vs open.}
\begin{itemize}
\item \textbf{Fixed in Cases 3--5:} production into the stock (once productive)
is \textbf{linear / immigration-type}, not birth--death exponential mean growth.
\item \textbf{Open / chosen by case:} release law --- continuous
($\varepsilon Q$) vs proportional full clear ($\delta Q^2/\Icell$); and whether
eclipse is present (Cases 2 and 5, or stacked on Case~3).
\item \textbf{Composability:} Case~5 is exactly Case~4 with Case~2's eclipse
block in front.
\item \textbf{Link to renewal BMVR:} each case induces kernels $S(a)$ and
$g(a)$; the convolution form of Result~\ref{res:new-vde} remains available when
ODE stage structure is replaced by age structure.
\end{itemize}

\paragraph{Takeaway.}
Start from classic BMVR (Case~1), add eclipse if delay matters (Case~2), then
--- when release timing and stock matter --- replace $p\Icell$ by an explicit
compartment $Q$ with either continuous export (Case~3) or proportional full
clear (Case~4). Add eclipse in front of full clear as Case~5 when both delay
and batch-like dump matter. Classic $p$ is recovered as a quasi-steady limit of
Case~3.

% ===========================================================================
\section{Self-excitatory release mechanisms}
\label{sec:self-excite}
% ===========================================================================

Cases 3--5 treat release as either continuous export or load-proportional full
clear. A further biological possibility is that release is partially
\textbf{self-excitatory}: one budding / dump event makes further release more
likely for a short time (Hawkes-type thinking), e.g.\ via cooperative Gag
nucleation, local membrane platforms, or ESCRT already recruited at a site.

This must be balanced against \textbf{depletion} of stock $X$ and finite
machinery: pure open-ended self-excitation is unlikely; a hybrid
``excitation + refill + depletion'' picture is more plausible than a classical
Hawkes process alone.

\paragraph{Two levels of BMVR.}
\begin{itemize}
\item \textbf{ODE BMVR:} finite system in
$(\Icell,Q\text{ or }m,\Vfree)$ (and maybe $E$, phase variables).
\item \textbf{Renewal BMVR:} $(\Icell,\Vfree)$ with kernels $S(a)$, $g(a)$ from
a single-cell process --- same skeleton as Result~\ref{res:new-vde}.
\end{itemize}
Renewal BMVR exists whenever the single-cell model defines an expected release
rate $g(a)$ and survival $S(a)$. ODE BMVR needs a low-dimensional Markov
state (or a moment closure).

\subsection{Menu of self-excitatory models}

\begin{center}
\renewcommand{\arraystretch}{1.2}
{\small
\begin{tabular}{@{}clp{0.28\textwidth}cc@{}}
\toprule
\# & Model & Excitation type & ODE BMVR? & Renewal? \\
\midrule
1 & Classical Hawkes & event history only & poor & yes (num.\ $g$) \\
2 & Hawkes $+$ stock $X$ & events $+$ depletion/refill & approx.\ / via 10 & \textbf{yes} \\
3 & Marked Hawkes & size-weighted & approx. & yes (num.) \\
4 & ETAS-style & size $+$ time kernel & approx. & yes (num.) \\
5 & Cox $+$ latent $Z$ & latent excitability & yes if $Z$ simple & yes \\
6 & ON/OFF telegraph & sustained high phase & \textbf{yes} & yes \\
7 & $\lambda\propto X^2$ & cooperative rate, no history & \textbf{yes} (Cases 4--5) & yes \\
8 & Nucleation biophysics & assembly cascade & hard & yes (num.) \\
9 & Age since last release & renewal excitation & yes (stages) & yes \\
10 & Jumping rate $r(t)$ & shot-noise intensity & \textbf{yes} (mean-field) & yes \\
11 & Spatial Hawkes & local membrane & average first & after coarse-grain \\
12 & Full clear (Cases 4--5) & instant total cascade & \textbf{yes} & yes \\
\bottomrule
\end{tabular}}
\end{center}

\paragraph{How they nest.}
Model 2 (Hawkes $+$ $X$) is the general event-history form. With
$\phi\equiv 0$ and $f(x)=\varepsilon x$ it recovers continuous export
(Case~3); with $\phi\equiv 0$, $f(x)=\delta x^2$ and full clear it recovers
model 7 / Cases 4--5; with exponential $\phi(t)=\eta e^{-\beta t}$ it is
equivalent to model 10 (Markov intensity). Model 6 is a coarse two-level
approximation to high vs low release rate without continuous $r$.

The four models best suited to our build are detailed next: \textbf{7},
\textbf{6}, \textbf{10}, and \textbf{2}.

% ---------------------------------------------------------------------------
\subsection{Model 7 --- Quadratic / cooperative intensity ($\lambda\propto X^2$)}
% ---------------------------------------------------------------------------

\paragraph{Idea.}
Release is a state-dependent Poisson process. Given load $X=x$,
\[
\lambda(x)=\delta x^2
\quad\text{(or }\delta x(x-1)\text{)}.
\]
No explicit memory of past event times --- only current load.
``Self-excitation'' means \textbf{cooperativity}: more material present makes
release disproportionately more likely (nucleation / lattice / local Gag
density). This is the single-cell version of the mean-field dump rate in
Cases~4--5.

\paragraph{Single-cell transitions (productive cell).}
\begin{center}
\renewcommand{\arraystretch}{1.15}
\begin{tabular}{@{}lll@{}}
\toprule
Event & Effect & Rate \\
\midrule
Production & $X\to X+1$ & $\nu$ (linear / immigration) \\
Optional degradation & $X\to X-1$ & $\mu X$ \\
Dump (full clear) & $X\to 0$, release $+X$ & $\delta X^2$ \\
Cell death & stop & $d_{\Icell}$ \\
\bottomrule
\end{tabular}
\end{center}

\paragraph{Population BMVR.}
ODE form is already Cases~4--5:
\[
\Vfree'
=\delta\frac{Q^2}{\Icell}-c\Vfree,
\qquad
Q'
=\nu\Icell-\delta\frac{Q^2}{\Icell}-d_{\Icell}Q
\]
(homogeneous closure $\E[X^2]\approx m^2$, $m=Q/\Icell$). Renewal form:
$g(a)=\delta\E[X_a^2\mathbf 1_{\{\mathrm{alive}\}}]$, $S(a)=e^{-d_{\Icell}a}$,
then the usual convolutions.

\paragraph{Pros / cons.}
ODE and renewal both available; interpretable cooperativity without Hawkes;
homogeneous $Q^2/\Icell$ ignores load variance; full clear is an extreme
(unit budding at rate $\varepsilon X$ is milder).

\paragraph{Use when.}
Default bulk ODE ``excitatory'' model; continuum of cooperative nucleation if
event history is not required.

% ---------------------------------------------------------------------------
\subsection{Model 6 --- Two-state telegraph / bursting release (ON--OFF)}
% ---------------------------------------------------------------------------

\paragraph{Idea.}
The cell (or a membrane site) switches between OFF (little or no release) and
ON (release at a high rate). Self-excitation is \textbf{phase-based}: once ON,
many releases occur until OFF --- a burst, not a Hawkes sum over every past
event. Optional: OFF$\to$ON rate increases with $X$ or after a recent release.

\paragraph{Single-cell sketch.}
State $(M,X)$ with $M\in\{\mathrm{OFF},\mathrm{ON}\}$. Production at rate
$\nu$ (always or only ON); OFF$\to$ON at rate $\alpha$ (or
$\alpha_0+\alpha_1 X$); ON$\to$OFF at rate $\beta$; release while ON at rate
$\varepsilon X$ or constant $\rho$; cell death $d_{\Icell}$.

\paragraph{Sample path.}
OFF: stock may build, little free virus; ON: rapid release, stock may fall;
OFF again: rebuild. Intermittent surges without requiring full clear.

\paragraph{Population ODE BMVR (simplest --- constant $\rho$ while ON).}
\begin{empheq}[box=\fbox]{align}
\ddt{I_{\mathrm{OFF}}}
&= \gamma T\,\Vfree + \beta I_{\mathrm{ON}}
- (\alpha+d_{\Icell})I_{\mathrm{OFF}},
\label{eq:m6-off}\\[0.25em]
\ddt{I_{\mathrm{ON}}}
&= \alpha I_{\mathrm{OFF}}
- (\beta+d_{\Icell})I_{\mathrm{ON}},
\label{eq:m6-on}\\[0.25em]
\ddt{\Vfree}
&= \rho\,I_{\mathrm{ON}} - c\,\Vfree.
\label{eq:m6-V}
\end{empheq}
(Eclipse: feed $E$ by incidence, convert into $I_{\mathrm{OFF}}$.) Quasi-steady
switching yields effective budding
\[
p_{\mathrm{eff}}
=\rho\cdot\frac{\alpha}{\alpha+\beta},
\]
recovering classic BMVR (Case~1) as a duty-cycle limit.

\paragraph{With stock.}
Track $Q$ with production $\nu(I_{\mathrm{ON}}+I_{\mathrm{OFF}})$ (or only ON)
and export $\varepsilon Q$ restricted to ON cells (requires a closure for how
load is shared). Renewal BMVR: single-cell $(M_a,X_a)$ defines $g$ and $S$.

\paragraph{Pros / cons.}
Among excitatory models, one of the \textbf{best for writable ODEs}; matches
intermittent release; links to Tat ON/OFF intuition; extra parameters
$(\alpha,\beta,\rho)$; not true Hawkes (no decaying event kernel).

\paragraph{Use when.}
Clustered release in time with clear ODEs and a transparent limit to constant
$p$.

% ---------------------------------------------------------------------------
\subsection{Model 10 --- Jumping continuous release rate $r(t)$}
% ---------------------------------------------------------------------------

\paragraph{Idea.}
Keep a continuous release-rate (intensity) process $r(t)\ge 0$: jumps up on
excitation events, decays toward a baseline between jumps. Export flux is
$r(t)$ or $r(t)X(t)$. This is the \textbf{Markovian form of Hawkes} when the
excitation kernel is exponential (shot-noise representation).

\paragraph{Single-cell (exponential kernel).}
\begin{align}
\mathrm{d}r
&= -\beta\bigl(r-r_0\bigr)\,\mathrm{d}t + \eta\,\mathrm{d}N_t,
\label{eq:m10-r}\\
\lambda(t)
&= r(t)
\quad\text{or}\quad
r(t)\,X_t,
\qquad
N\text{ has intensity }\lambda,
\label{eq:m10-lam}
\end{align}
with $X$ refilled at rate $\nu$ and decremented on each release. Each event
adds $\eta$ to $r$; memory fades at rate $\beta$.

\paragraph{Population ODE BMVR (mean-field sketch).}
Track $\Icell$, total load $Q$, total excitation $R=\sum r_i$ (mean
$\bar r=R/\Icell$). With flux $\approx \bar r\,Q$ and jump input
$\propto$ event rate $\approx\bar r\,Q$:
\begin{empheq}[box=\fbox]{align}
\ddt{\Icell}
&= \gamma T\,\Vfree - d_{\Icell}\,\Icell,
\label{eq:m10-I}\\[0.25em]
\ddt{Q}
&= \nu\Icell - \bar r\,Q - d_{\Icell}Q,
\label{eq:m10-Q}\\[0.25em]
\ddt{R}
&= -\beta\bigl(R-r_0\Icell\bigr) + \eta\,(\bar r\,Q)
+\text{(cell birth/death corrections)},
\label{eq:m10-R}\\[0.25em]
\ddt{\Vfree}
&= \bar r\,Q - c\,\Vfree,
\label{eq:m10-V}
\end{empheq}
using closures $\E[rX]\approx\bar r\,m$. Renewal BMVR: $g(a)=\E[r_a X_a]$
(if flux $=rX$); kernels usually numerical.

\paragraph{Pros / cons.}
True self-excitation with fading memory; finite-dimensional intensity state
allows ODE BMVR; bridges Hawkes and bulk models; needs closures and extra
parameters $(\beta,\eta,r_0)$; hard to identify from viral load alone.

\paragraph{Use when.}
Hawkes-type memory is desired but host-level equations should remain ODEs.

% ---------------------------------------------------------------------------
\subsection{Model 2 --- Hawkes process with stock $X$}
% ---------------------------------------------------------------------------

\paragraph{Idea.}
Release times $(t_i)$ with intensity
\begin{equation}
\label{eq:m2-hawkes}
\lambda(t)
= f\bigl(X_t\bigr)
+ \sum_{t_i<t}\phi(t-t_i),
\end{equation}
$f(X)$ baseline from load (e.g.\ $f(x)=\varepsilon x$ or $\delta x^2$); $\phi\ge 0$
self-excitation kernel; on event, $X\leftarrow X-1$ (or batch / full clear);
between events, $X$ grows at rate $\nu$; cell dies at rate $d_{\Icell}$.
This is the most literal ``Hawkes for budding with biology.''

\paragraph{Special cases.}
\begin{center}
\renewcommand{\arraystretch}{1.15}
\begin{tabular}{@{}ll@{}}
\toprule
Choice & Recovers \\
\midrule
$\phi=0$, $f(x)=\varepsilon x$ & Case~3 continuous export \\
$\phi=0$, $f=\delta x^2$, jump $X\to 0$ & Model 7 / Cases 4--5 \\
$f=\mathrm{const}$, $\phi=\eta e^{-\beta t}$, ignore $X$ & classical Hawkes \\
$f=\varepsilon x$, $\phi=\eta e^{-\beta t}$ & \textbf{canonical hybrid} \\
exponential $\phi$ & equivalent to Model 10 \\
\bottomrule
\end{tabular}
\end{center}

\paragraph{Population BMVR.}
\begin{itemize}
\item \textbf{Renewal:} always well-defined ---
$g(a)=\E[\lambda(a)\mid\mathrm{alive}]$, $S$ from cell lifetime; same
convolutions as BDC. Kernels typically from simulation (or from Model 10 if
$\phi$ is exponential).
\item \textbf{ODE:} not without reduction --- infinite event history. Restrict
to exponential $\phi$ and use Model 10 mean-field, or apply moment closures.
\end{itemize}

\paragraph{Pros / cons.}
Direct Hawkes language; flexible short excitation window (e.g.\ ESCRT);
combines load, excitation, and depletion; analytics hard unless exponential
$\phi$; ODE BMVR only after reduction to Model 10.

\paragraph{Use when.}
Conceptual / single-cell / renewal theory with true self-excitation; for host
ODEs, specialise immediately to exponential $\phi$ and Model 10.

% ---------------------------------------------------------------------------
\subsection{Comparison of the four focal models}
% ---------------------------------------------------------------------------

\begin{center}
\renewcommand{\arraystretch}{1.25}
{\small
\begin{tabular}{@{}p{0.18\textwidth}p{0.18\textwidth}p{0.18\textwidth}p{0.18\textwidth}p{0.18\textwidth}@{}}
\toprule
 & \textbf{7} $\lambda\propto X^2$ & \textbf{6} ON/OFF & \textbf{10} jumping $r$ & \textbf{2} Hawkes$+X$ \\
\midrule
Carrier & load only & discrete phase & continuous $r$ & event history ($+$ load) \\
Memory & Markov in $X$ & Markov in phase & Markov in $r$ & non-Markov unless exp.\ $\phi$ \\
Clusters & high $X$ $\to$ fast dumps & long ON sojourn & high-$r$ window & aftershocks \\
ODE BMVR & yes (Cases 4--5) & yes (split $I$) & yes (mean-field $R$) & only via 10 \\
Renewal BMVR & yes & yes & yes & yes \\
$\to$ classic $p$ & QSS $\delta m^2$ & duty cycle & QSS $\E[rX]$ & as 10 if exp $\phi$ \\
\bottomrule
\end{tabular}}
\end{center}

\paragraph{Biological confidence (shared caveat).}
Cooperative assembly / nucleation (Gag lattice, scaffolds, ESCRT) makes
\textbf{short clustered} release plausible (\textbf{medium--high} as a
hypothesis). A classical Hawkes process with only event memory and no
depletion is a poor standalone cell model (\textbf{low--medium}). Hybrid
excitation $+$ stock depletion/refill is a reasonable working hypothesis; it is
\textbf{not} established as the dominant single-cell export law for HIV.

\paragraph{Practical ranking for this project.}
\begin{enumerate}
\item Writable ODEs inside current Cases 3--5: Model \textbf{7}.
\item Intermittent release, simplest new ODEs: Model \textbf{6}.
\item Self-excitation with fading memory $+$ ODEs: Model \textbf{10}.
\item Paper language ``Hawkes-like budding'': define Model \textbf{2},
implement ODEs as Model \textbf{10}, or renewal with numerical $g$.
\end{enumerate}

\paragraph{Takeaway.}
Self-excitation is compatible with BMVR whenever the single cell supplies
$(S,g)$ (renewal form always). Finite ODE BMVR needs excitation carried by a
few Markov states --- load $X$ (Model 7), ON/OFF (Model 6), or intensity $r$
(Model 10) --- not an unbounded event history, unless that history is reduced
to exponential memory (Hawkes $\leftrightarrow$ Model 10).

% ===========================================================================
\section{Partial release: intermediates between continuous export and full clear}
\label{sec:partial-release}
% ===========================================================================

Cases 3--5 sit at two extremes of \emph{how much} leaves the intracellular
compartment when release occurs:

\begin{center}
\renewcommand{\arraystretch}{1.2}
\begin{tabular}{@{}lll@{}}
\toprule
Extreme & Mechanism & In this document \\
\midrule
Continuous export & units leave steadily (flux $\varepsilon Q$) & Case 3 \\
Full clear & event empties the whole pool & Cases 4--5 \\
\bottomrule
\end{tabular}
\end{center}

A natural intermediate is a \textbf{release event} that removes only a
\textbf{random or fixed part} of the current load $X=n$: export $B$ units with
$0\le B\le n$ (or $1\le B\le n$), retain $n-B$. Self-excitation
(Section~\ref{sec:self-excite}) is optional and separate; here only the
\textbf{release-count law} $B\mid X=n$ is at issue.

\paragraph{Two ingredients.}
\begin{enumerate}
\item \textbf{When} an event fires --- rate $\lambda(X)$, e.g.\ constant
$\delta$, linear $\delta X$, or cooperative $\delta X^2$.
\item \textbf{How many} leave --- distribution of $B\mid X=n$.
\end{enumerate}

% ---------------------------------------------------------------------------
\subsection{Options for $B\mid X=n$}
% ---------------------------------------------------------------------------

\subsubsection*{A. Boolean / independent retention (recommended default intermediate)}

When an event fires at load $n$, each of the $n$ units is released independently
with probability $q$ and retained with probability $1-q$:
\begin{equation}
\label{eq:boolean-B}
B\mid X=n \;\sim\; \mathrm{Binomial}(n,q),
\qquad
X \leftarrow n-B \;\sim\; \mathrm{Binomial}(n,1-q).
\end{equation}
\begin{center}
\renewcommand{\arraystretch}{1.15}
\begin{tabular}{@{}ll@{}}
\toprule
$q$ & Behaviour \\
\midrule
$q\to 0$ & tiny fractional clear \\
$q\in(0,1)$ & true intermediate \\
$q=1$ & full clear (Cases 4--5 style) \\
\bottomrule
\end{tabular}
\end{center}
Mean release given $n$: $\E[B\mid n]=qn$. Mechanistic reading: a release event
``hits'' the pool; each package succeeds independently with probability $q$.
Generating functions stay friendly (product structure over units).

\paragraph{Event rate still free.} Examples: $\lambda=\delta$ (per cell);
$\lambda=\delta X$; $\lambda=\delta X^2$.

\paragraph{Mean-field BMVR (boolean, $\lambda=\delta X$).}
Population event structure as in Case~4, but only fraction $q$ of load leaves:
\begin{empheq}[box=\fbox]{align}
\ddt{\Icell}
&= \gamma T\,\Vfree - d_{\Icell}\,\Icell,
\label{eq:bool-I}\\[0.25em]
\ddt{Q}
&= \nu\Icell - \delta q\,\frac{Q^2}{\Icell} - d_{\Icell}Q,
\label{eq:bool-Q}\\[0.25em]
\ddt{\Vfree}
&= \delta q\,\frac{Q^2}{\Icell} - c\,\Vfree.
\label{eq:bool-V}
\end{empheq}
Thus $q=1$ recovers Case~4; $q\in(0,1)$ is partial clear with the same clock.
If instead $\lambda=\delta$ (constant per cell), mean export is
$\delta q\,Q$ and the bulk ODEs collapse to Case~3 with
$\varepsilon_{\mathrm{eff}}=\delta q$ --- pulses remain visible in single-cell
paths and higher moments, not in the first-moment ODE.

\subsubsection*{B. Fixed fraction}

On event, release $B=\lceil fn\rceil$ or $\mathrm{round}(fn)$ with fixed
$f\in(0,1]$. Same mean as boolean with $q=f$, \textbf{less} stochasticity.
Bulk mean-field matches boolean at equal mean fraction; single-cell burst-size
variance is smaller.

\subsubsection*{C. Random fraction (continuous thinning)}

Draw $U\sim\mathrm{Beta}(a,b)$ (or $\mathrm{Uniform}(0,1)$); set
$B=\mathrm{round}(Un)$ or thin continuously. More variable burst sizes than
fixed $f$; boolean is the discrete ``each particle flips a coin'' analogue of
random thinning.

\subsubsection*{D. Fixed batch size (capped)}

$B=\min(k,n)$ for package size $k\ge 1$. Special cases: $k=1$ is unit budding
on an event clock (Case~3 in the frequent-event limit); $k$ large approaches
full clear when $n$ is not huge.

\subsubsection*{E. Size-biased / prefer large}

$P(B=j\mid n)\propto j$, or dump probability depending strongly on $n$. Less
natural if units are identical; more natural if one site releases a random
assembled blob.

\subsubsection*{F. Threshold then partial}

Event only if $X\ge x_*$, then apply boolean $q$ or fixed $f$. Combines
accumulation with a trigger (``surge when enough material'').

\subsubsection*{G. Short high-export window}

On trigger, export at high $\varepsilon$ for a short time, then stop. Partial
clear in expectation without an explicit law for $B$; related to ON/OFF
(Model~6) with brief ON.

% ---------------------------------------------------------------------------
\subsection{Single-cell boolean process (compact)}
% ---------------------------------------------------------------------------

\begin{center}
\renewcommand{\arraystretch}{1.2}
\begin{tabular}{@{}lll@{}}
\toprule
Event & Rate & Jump \\
\midrule
Produce & $\nu$ & $X\to X+1$ \\
Release event & $\lambda(X)$ (e.g.\ $\delta X$ or $\delta X^2$) &
  $B\sim\mathrm{Bin}(n,q)$; $X\to n-B$; free virus $+B$ \\
Cell death & $d_{\Icell}$ & stop \\
\bottomrule
\end{tabular}
\end{center}

\paragraph{Nested limits.}
\begin{center}
\renewcommand{\arraystretch}{1.15}
\begin{tabular}{@{}ll@{}}
\toprule
Choice & Recovers \\
\midrule
$q=1$, $\lambda=\delta X^2$ & full clear, cooperative (Case 4-like) \\
$q=1$, $\lambda=\delta X$ & full clear, linear hazard \\
$0<q<1$, $\lambda=\delta X$ & \textbf{main partial-clear intermediate} \\
always $B=1$ when fire, $\lambda=\varepsilon X$ & unit budding (Case 3 discrete) \\
\bottomrule
\end{tabular}
\end{center}

% ---------------------------------------------------------------------------
\subsection{BMVR feasibility}
% ---------------------------------------------------------------------------

\begin{center}
\renewcommand{\arraystretch}{1.2}
\begin{tabular}{@{}p{0.42\textwidth}cc@{}}
\toprule
Setup & ODE BMVR & Renewal BMVR \\
\midrule
Boolean $q$, $\lambda=\delta$ (const) & yes ($\approx$ Case 3, $\varepsilon=\delta q$) & yes \\
Boolean $q$, $\lambda=\delta X$ & yes \eqref{eq:bool-I}--\eqref{eq:bool-V} & yes \\
Boolean $q$, $\lambda=\delta X^2$ & yes (Case 4 factor $q$, retain $1-q$) & yes \\
General law $B\mid n$ & yes if $\E[B\mid n]$ closes under a moment assumption & yes always via $g(a)$ \\
\bottomrule
\end{tabular}
\end{center}

Renewal BMVR: single-cell process defines
$g(a)=\E[\text{release rate at age }a]$ and survival $S(a)$; population
equations as in Result~\ref{res:new-vde}.

% ---------------------------------------------------------------------------
\subsection{Where partial clear sits relative to continuous export}
% ---------------------------------------------------------------------------

\begin{center}
\renewcommand{\arraystretch}{1.2}
\begin{tabular}{@{}lll@{}}
\toprule
Model & Single-cell path & Mean bulk ODE \\
\midrule
Case 3 ($\varepsilon Q$) & smooth drip & linear export \\
Boolean, $\lambda=\delta$, any $q$ & rare binomial pulses & $\approx$ linear $\delta q\,Q$ \\
Boolean, $\lambda=\delta X$, $q\in(0,1)$ & pulses of mean size $qm$ & quadratic, strength $\delta q$ \\
Case 4 ($q=1$) & full empty pulses & quadratic full clear \\
\bottomrule
\end{tabular}
\end{center}

Partial vs full is not always visible in mean ODEs if the event rate is
constant (then only single-cell burst sizes and higher moments differ). It
\textbf{is} visible in bulk dynamics when the event rate depends on load
(e.g.\ $\lambda=\delta X$), and always in single-cell release-count
distributions.

% ---------------------------------------------------------------------------
\subsection{Recommendation}
% ---------------------------------------------------------------------------

\textbf{Primary intermediate:} boolean (binomial) thinning with event rate
$\lambda=\delta X$ (or $\delta X^2$) and release probability $q\in(0,1]$.
One transparent parameter $q$ between full clear and mild partial clear;
mechanistic; nested in Cases 3--4 at mean-field level; easy to simulate and to
write BMVR ODEs.

\textbf{Secondary:} fixed fraction $f$ if less burst-size noise than binomial is
desired.

\textbf{Later / data-driven:} Beta random fraction, capped batch $k$, threshold
rules, or short high-export windows --- unless measurements demand them.

\paragraph{Takeaway.}
Between continuous budding and full dumping, let a release event remove a
\textbf{distributed part} of the pool. The clean default is independent
boolean release of each unit with probability $q$; with a load-dependent event
clock this is Case~4 scaled by $q$, and with a constant clock it is mean-equivalent
to Case~3 while still producing pulsed single-cell paths.

% ===========================================================================
\section{What is publishable, and how it fits together}
% ===========================================================================

\subsection{Core claims suitable for a paper}

\begin{enumerate}
\item \textbf{Closed-form single-cell theory} for the BDC, including simplified
moments, corrected non-fixation probability, and a complete proof of the
geometric burst-size law.
\item \textbf{Identity:} the burst-size distribution equals the quasi-stationary
distribution (geometric with ratio $1/a$).
\item \textbf{Renewal BMVR} with closed-form kernels $\Ihat$ and $\delta K$,
exact effective parameters $p_{\rm eff}(r)$ and $d_{\Icell,{\rm eff}}(r)$, and
invariant $R_0$.
\item \textbf{Growth-phase dependence of~$p$} as a concrete, directional
prediction for BMVR fitting practice.
\item \textbf{Identifiability:} three intracellular rates cannot be recovered
from two BMVR parameters alone; three rescuing observables are named.
\item \textbf{Flooding advantage} (Result~\ref{res:flood}): the exact condition
under which bursting beats matched-$R_0$ budding for early establishment.
\end{enumerate}

\subsection{Suggested paper structure}

Target: \textit{Journal of Theoretical Biology} or
\textit{Bulletin of Mathematical Biology}.

\begin{enumerate}
\item Introduction --- BMVR's constant-release assumption; single-cell evidence;
what is and is not new (renewal framework standard; closed-form kernels and
consequences new).
\item The intracellular process --- BDC, closed forms, QSD/burst-size identity.
\item Renewal reformulation --- kernels and equations.
\item Exact reduction --- $p_{\rm eff}$, $d_{\Icell,{\rm eff}}$, characteristic
equation, $R_0$ unchanged, $p$ growth-phase dependent.
\item Identifiability --- negative result and rescuing observables.
\item The flooding advantage --- main theorem and the condition
$1/\delta<1/\beta+1/\mu$.
\item Application --- single-cell release data; lytic vs.\ budding caveat.
\item Discussion --- limits of linear intracellular growth; two-type/eclipse
extension as future work (or fold in if completed).
\end{enumerate}

If the two-type kernels are finished, a natural split is: (i)~one-type theory as
above; (ii)~two-type kernel with a mechanistic eclipse phase.

\subsection{What this summary deliberately omits}

\begin{itemize}
\item Full errata catalogue for thesis Chapters~1, 3--6 (recorded in the working
notes).
\item Open technical tasks (literature check, conceptual QSD proof, two-type
moments, intermediate release models, nonlinear intracellular growth).
\item The $N=2$ immediate-transfer chain model (interesting, not yet at the same
state of completion).
\item Detailed hypergeometric expressions for the Laplace transforms (available
in the notes; not needed for the main story).
\end{itemize}

% ===========================================================================
\section{Quick-reference formula sheet}
% ===========================================================================

\begin{center}
\renewcommand{\arraystretch}{1.4}
\begin{tabular}{@{}lll@{}}
\toprule
Quantity & Closed form & Role \\
\midrule
$a,b$ & roots of the characteristic quadratic & structure \\
$\Aa,\Bb$ & $a-1$, $1-b$ & simplify \\
$I(t)$ & $(a\Bb+b\Aa w)/(\Bb+\Aa w)$ & no-burst survival \\
$D(t)$ & $ab(w-1)/(aw-b)$ & internal clearance \\
$\Ihat(t)$ & $(a-b)^2 w/[(\Bb+\Aa w)(aw-b)]$ & productive survival \\
$J(t)$ & $(a-b)^2 w/(\Bb+\Aa w)^2$ & mean load \\
$V(t)$ & $(1-I)[1+(\beta/\delta)(1-I)]$ & mean release \\
$V_\infty$ & $a(1-b)/(a-1)$ & total mean output \\
$\varphi(t)$ & $\delta J(t)$ & burst-time density \\
$\Prob{\Kb=k}$ & $(\delta/\beta)a^{-k}$ & burst size \\
QSD & $(a-1)a^{-n}$ & conditional load \\
$\E[T_{\rm prod}]$ & $\beta^{-1}\log(a/(a-1))$ & productive lifetime \\
$R_0$ & $\gamma T V_\infty/c$ & threshold (unchanged) \\
$p_{\rm eff}(r)$ & $\delta\Lap{K}/\Lap{\Ihat}$ & effective release \\
Flooding & $1/\delta < 1/\beta+1/\mu$ & when bursting wins \\
\midrule
\multicolumn{3}{@{}l@{}}{\textit{New viral dynamics (Result~\ref{res:new-vde})}} \\
$\Icell(t)$ & $\Icell_0\Ihat(t)+(\gamma T\,\Vfree*\Ihat)(t)$ & infected cells \\
$\Vfree'(t)$ & $\Icell_0 g(t)+(\gamma T\,\Vfree*g)(t)-c\Vfree$ & free virions \\
$g(a)$ & $\delta K(a)=V'(a)$ & release kernel (BDC) \\
\midrule
\multicolumn{3}{@{}l@{}}{\textit{Reset catastrophe (Section~\ref{sec:reset-bmvr})}} \\
$S(a)$ & $e^{-d_{\Icell}a}$ & cell survival \\
$g(a)$ & $\delta e^{-d_{\Icell}a}\E[X_a^2]$ & multi-pulse release \\
BMVR & same convolutions as BDC & kernels only change \\
$p_\infty$ & $\delta\E_\pi[X^2]$ & mature-cell limit \\
\midrule
\multicolumn{3}{@{}l@{}}{\textit{HIV coarse grain (Section~\ref{sec:hiv-mean-growth})}} \\
Intracellular mean & eclipse $\to$ $m'\approx\nu-\mu m$ & not $m'=rm$ \\
Tat intermediate & rate $\le$ late ON ceiling & no systematic overshoot \\
\bottomrule
\end{tabular}
\end{center}

\vspace{1.5em}
\begin{center}
\textit{End of summary. Full derivations, numerical verification tables,\\
and open tasks live in the BMVR extension working notes.}
\end{center}

\end{document}
